\documentclass[12pt,a4paper,UTF8]{ctexart}
\usepackage[a4paper,top=22mm,bottom=22mm,left=22mm,right=22mm,headheight=15pt]{geometry}
\usepackage{amsmath,amssymb,mathtools,bm}
\usepackage{graphicx}
\usepackage{xcolor}
\usepackage{booktabs,longtable,array}
\usepackage{enumitem}
\usepackage{fancyhdr}
\usepackage{caption}
\usepackage{float}
\usepackage{microtype}
\usepackage{setspace}
\usepackage[hidelinks,unicode]{hyperref}
\usepackage[most]{tcolorbox}

\definecolor{ContestBlue}{HTML}{17365D}
\definecolor{SoftBlue}{HTML}{EEF4FA}
\definecolor{SoftGray}{HTML}{F5F6F7}
\hypersetup{pdftitle={第100套},pdfauthor={物理竞赛教学资料}}

\setstretch{1.24}
\setlength{\parindent}{2em}
\setlength{\parskip}{0.35em}
\setlist{itemsep=0.35em,topsep=0.4em}
\captionsetup{font=small,labelfont=bf,justification=centering}
\setcounter{secnumdepth}{0}

\pagestyle{fancy}
\fancyhf{}
\fancyhead[L]{\small 第100套}
\fancyfoot[C]{\thepage}
\renewcommand{\headrulewidth}{0.4pt}

\ctexset{
  section={format=\Large\bfseries\color{ContestBlue},beforeskip=1.2em,afterskip=0.7em},
  subsection={format=\large\bfseries\color{ContestBlue},beforeskip=1.0em,afterskip=0.5em},
  subsubsection={format=\normalsize\bfseries,beforeskip=0.8em,afterskip=0.35em}
}

\makeatletter
\newcommand{\tightlist}{\setlength{\itemsep}{0pt}\setlength{\parskip}{0pt}}
\newcommand{\maxwidth}{\ifdim\Gin@nat@width>\linewidth\linewidth\else\Gin@nat@width\fi}
\newcommand{\maxheight}{\ifdim\Gin@nat@height>0.72\textheight 0.72\textheight\else\Gin@nat@height\fi}
\newcommand{\pandocbounded}[1]{\begingroup\setkeys{Gin}{width=\maxwidth,height=\maxheight,keepaspectratio}#1\endgroup}
\makeatother

\begin{document}

\begin{titlepage}
  \centering
  \vspace*{58mm}
  {\Huge\bfseries\color{ContestBlue} 第100套\par}
  \vfill
  \vspace*{28mm}
\end{titlepage}


\section{第 1 题}\label{ux7b2c-1-ux9898}

半径为 \(R\)
的半球形山坡表面光滑，坡顶有一小洞。把高尔夫球视为质点，球最初位于半球外表面上，球心位置对应的半径与竖直向上的半径夹角为
\(\theta\)。现沿球面切线、指向坡顶的方向给球初速度 \(v\)，重力加速度为
\(g\)。

\begin{enumerate}
\def\labelenumi{\arabic{enumi}.}
\tightlist
\item
  若球始终贴着半球表面上滑，并在到达坡顶时速度恰好减为零，求初速度
  \(v\)。
\item
  求第（1）问方案可行时 \(\theta\) 的范围。
\item
  当第（1）问方案不可行时，应怎样使球进入坡顶小洞？求新方案所需的初速度
  \(v\)。
\item
  求第（3）问方案可行时 \(\theta\) 的范围。
\end{enumerate}

\begin{figure}[H]
\centering
\includegraphics[width=0.698\linewidth,height=\textheight,keepaspectratio,alt={图 1-1　小球从半球外表面沿切线方向射向坡顶。}]{../figures/question-01.pdf}
\caption*{图 1-1　小球从半球外表面沿切线方向射向坡顶。}
\end{figure}

\subsection{一、题型定位}\label{ux4e00ux9898ux578bux5b9aux4f4d}

小球先沿外侧圆弧运动，再沿脱离点的切线作抛射。方案能否成立由支持力非负决定；临界能量还须检查是否能在有限时间到达坡顶。

\subsection{二、知识点}\label{ux4e8cux77e5ux8bc6ux70b9}

\begin{itemize}
\tightlist
\item
  光滑固定曲面上运动的机械能守恒。
\item
  外圆弧径向方程和支持力约束 \(N\geq0\)。
\item
  接触失效条件 \(N=0\)，以及``算出的 \(N<0\)''表示物体已经离开曲面。
\item
  斜抛运动的分方向方程。
\item
  利用初始切线方向建立速度分量。
\item
  轨迹可行性的端点、方向和区间检查。
\end{itemize}

\subsection{三、物理图像}\label{ux4e09ux7269ux7406ux56feux50cf}

沿球面上滑时，速度越大，指向球心的向心加速度越大。重力只能提供
\(mg\cos\varphi\)
的向心分量，而曲面对小球的支持力指向球心外侧，不能帮助提供向心力。因此，小球若在较低处速度过大，维持圆周轨迹所需的支持力会变成负值；真实运动只能是离开球面。

后两问中，小球从初始点沿切线飞出。此时半球只负责确定发射点和发射方向，之后的轨迹完全由重力决定。目标条件是抛物线经过坡顶，而不是``到坡顶速度为零''。

\subsection{四、解题入口}\label{ux56dbux89e3ux9898ux5165ux53e3}

把运动中任意位置的圆心角记为 \(\varphi\)，初始时
\(\varphi=\theta\)，坡顶处
\(\varphi=0\)。先利用``到顶速度为零''反推初速度，再把同一能量关系代入径向方程求
\(N(\varphi)\)。接触条件要在整个区间 \(0\leq\varphi\leq\theta\) 上检查。

若接触条件失败，则在以半球圆心为原点的直角坐标系中写初始点、坡顶和切向速度分量，消去飞行时间。

\subsection{五、解析}\label{ux4e94ux89e3ux6790}

图 1-2 左侧画出小球贴外球面上滑时沿径向的重力分量 \(mg\cos\varphi\)
与支持力
\(N\)；右侧画出小球从初始点沿切线离面后进入坡顶小洞的抛物线，并标出水平、竖直几何位移。第（1）、（2）步使用左图同时检查能量与接触条件，第（3）、（4）步使用右图建立抛体方程。

\begin{figure}[H]
\centering
\includegraphics[width=0.9\linewidth,height=\textheight,keepaspectratio,alt={图 1-2　贴面上滑时的径向受力，以及切线抛射方案的几何关系。}]{../figures/solution-01.pdf}
\caption*{图 1-2　贴面上滑时的径向受力，以及切线抛射方案的几何关系。}
\end{figure}

\textbf{（1）贴着球面上滑并在坡顶停下}

按图 1-2
左侧，机械能方程只能确定各位置的速率；曲面能否维持接触还要由径向方程求
\(N\)。因此本步先用到顶速度为零确定初速度，下一步再逐点检查支持力。

任意位置的高度相对圆心水平线为 \(R\cos\varphi\)。从初始位置
\(\varphi=\theta\) 上升到任意位置 \(\varphi\)，机械能守恒给出

\[\frac12mv^2(\varphi)+mgR\cos\varphi
=\frac12mv^2+mgR\cos\theta. \tag{1}\]

题目要求小球到达坡顶时速度恰好为零。令
\(\varphi=0\)、\(v(0)=0\)，式（1）变成

\[mgR=\frac12mv^2+mgR\cos\theta.\]

移项并约去质量 \(m\)：

\[\frac12v^2=gR(1-\cos\theta).\]

因此

\[\boxed{v=\sqrt{2gR(1-\cos\theta)}}. \tag{2}\]

这一步只保证能量上``恰好到顶''，尚未保证运动过程中始终贴着球面。

\textbf{（2）贴面方案的可行范围}

图 1-2 左侧中，指向球心的重力分量为
\(mg\cos\varphi\)，而外球面对小球的支持力背离球心，所以径向方程必须写成
\(mg\cos\varphi-N=mv^2/R\)。若算得
\(N<0\)，就意味着曲面必须``拉住''小球，接触方案失效。

由式（1）和式（2），任意位置的速率为

\[\begin{aligned}
v^2(\varphi)
&=v^2-2gR(\cos\varphi-\cos\theta)\\
&=2gR(1-\cos\theta)-2gR(\cos\varphi-\cos\theta)\\
&=2gR(1-\cos\varphi).
\end{aligned} \tag{3}\]

沿指向球心的径向取正方向。重力的径向分量为 \(mg\cos\varphi\)，支持力
\(N\) 指向球心外侧，因此

\[mg\cos\varphi-N=\frac{mv^2(\varphi)}{R}. \tag{4}\]

代入式（3）：

\[\begin{aligned}
N
&=mg\cos\varphi-\frac{m}{R}\,2gR(1-\cos\varphi)\\
&=mg(3\cos\varphi-2).
\end{aligned} \tag{5}\]

在上滑过程中，\(\varphi\) 从 \(\theta\) 逐渐减小到 \(0\)，所以
\(\cos\varphi\) 逐渐增大。式（5）中的支持力最小值出现在初始点
\(\varphi=\theta\)。要使全过程 \(N\geq0\)，只需且必须满足

\[3\cos\theta-2\geq0.\]

故

\[\boxed{0\leq\theta\leq\arccos\frac23}. \tag{6}\]

当 \(\theta=\arccos(2/3)\) 时，初始瞬间 \(N=0\)，随后
\(N>0\)，属于贴面方案的临界情形。

这里还可以把支持力随位置的变化单独核验一次。由式（5）直接求导：

\[\frac{\mathrm dN}{\mathrm d\varphi}=-3mg\sin\varphi. \tag{6a}\]

在 \(0<\varphi<\theta<\pi/2\)
内，\(\mathrm dN/\mathrm d\varphi<0\)。这句话的自变量是
\(\varphi\)，而小球实际运动时 \(\varphi\) 在减小，因此沿运动方向 \(N\)
反而逐渐增大。初态支持力为

\[N_0=mg(3\cos\theta-2),\]

坡顶支持力为

\[N_{\rm top}=mg.\]

所以只检查 \(N_0\geq0\) 的结论不是经验判断，而是由单调性严格保证的。若
\(N_0<0\)，数学方程要求曲面``向内拉''小球；普通接触面不能提供这种力，因此能量方案虽有形式解，却没有物理实现。

\textbf{（3）不能贴面时的切线抛射方案}

转看图 1-2
右侧。离面瞬间的速度沿球面切线，初始点到洞口的水平位移与竖直位移均由半径
\(R\) 和角度 \(\theta\)
给出；分别写两个分量方程并消去飞行时间，才能确定所需初速度。

当

\[\theta>\arccos\frac23 \tag{7}\]

时，若仍使用式（2）的初速度，初始支持力将为负，故小球不可能按设想贴着球面运动。可改为给小球一个合适的切向速度，使它从初始点立即离开半球，沿抛物线落入坡顶小洞。

取水平向右为 \(x\) 轴正方向、竖直向上为 \(y\)
轴正方向。初始点和坡顶坐标分别为

\[P(R\sin\theta,R\cos\theta),\qquad H(0,R).\]

指向坡顶的切线单位向量为

\[\boldsymbol e_t=(-\cos\theta,\sin\theta),\]

所以初速度分量为

\[v_x=-v\cos\theta,\qquad v_y=v\sin\theta. \tag{8}\]

设飞行时间为 \(t\)。水平方向位移满足

\[-R\sin\theta=-v\cos\theta\,t,\]

从而

\[t=\frac{R\tan\theta}{v}. \tag{9}\]

竖直方向位移为 \(R(1-\cos\theta)\)，故

\[R(1-\cos\theta)=v\sin\theta\,t-\frac12gt^2. \tag{10}\]

把式（9）代入式（10）：

\[R(1-\cos\theta)
=R\frac{\sin^2\theta}{\cos\theta}
-\frac{gR^2\tan^2\theta}{2v^2}. \tag{11}\]

将重力项移到左边，并通分化简：

\[\begin{aligned}
\frac{gR^2\tan^2\theta}{2v^2}
&=R\left[\frac{\sin^2\theta}{\cos\theta}-(1-\cos\theta)\right]\\
&=R\left[\frac{1-\cos^2\theta-\cos\theta+\cos^2\theta}
{\cos\theta}\right]\\
&=R\frac{1-\cos\theta}{\cos\theta}.
\end{aligned}\]

为避免在最后一步漏掉三角恒等式，可把求 \(v^2\)
的过程再拆成两步。先由上式得到

\[v^2=\frac{gR\tan^2\theta\cos\theta}{2(1-\cos\theta)}.\]

再利用

\[\tan^2\theta\cos\theta
=\frac{\sin^2\theta}{\cos\theta}
=\frac{(1-\cos\theta)(1+\cos\theta)}{\cos\theta},\]

约去公共因子 \(1-\cos\theta\)，才得到式（12）。这一约分要求
\(\theta\ne0\)；本分支本来就满足
\(\theta>\arccos(2/3)\)，因而没有额外丢解。

于是

\[v^2
=\frac{gR}{2}\frac{\tan^2\theta\cos\theta}{1-\cos\theta}
=\frac{gR(1+\cos\theta)}{2\cos\theta}. \tag{12}\]

因此新方案所需初速度为

\[\boxed{v=\sqrt{\frac{gR(1+\cos\theta)}{2\cos\theta}}
=\cos\frac\theta2\sqrt{\frac{gR}{\cos\theta}}}. \tag{13}\]

还可以用抛物线轨迹方程作独立核验。令 \(X\)
表示从发射点向左的水平位移、\(Y\) 表示相对发射点向上的竖直位移，则

\[X=v\cos\theta\,t,\]

\[Y=v\sin\theta\,t-\frac12gt^2.\]

由第一式消去 \(t\)：

\[t=\frac{X}{v\cos\theta}.\]

代入第二式得

\[Y=X\tan\theta-
\frac{gX^2}{2v^2\cos^2\theta}.\]

坡顶小洞相对发射点的坐标为

\[X_H=R\sin\theta,\]

\[Y_H=R(1-\cos\theta).\]

将 \(X_H,Y_H\) 代入轨迹方程：

\[R(1-\cos\theta)
=\frac{R\sin^2\theta}{\cos\theta}
-\frac{gR^2\sin^2\theta}{2v^2\cos^2\theta},\]

与式（11）完全一致，因而再次得到式（13）。这种写法把``命中小洞''直接表现为目标点在抛物线上。

到达小洞时的竖直速度也能帮助检查轨迹形状：

\[v_{y,H}=v\sin\theta-gt.\]

利用式（9）和式（12），两边乘以 \(v\)：

\[\begin{aligned}
v\,v_{y,H}
&=v^2\sin\theta-gR\tan\theta\\
&=\frac{gR\sin\theta}{2\cos\theta}
(1+\cos\theta)-\frac{gR\sin\theta}{\cos\theta}\\
&=-\frac{gR\sin\theta(1-\cos\theta)}{2\cos\theta}<0.
\end{aligned}\]

因此球到达坡顶小洞时已经处在抛物线下降段，而不是恰好在最高点。这个结果与图中的轨迹形状一致。

\textbf{（4）切线抛射方案的可行范围}

式（13）要求 \(\cos\theta>0\)，即
\(\theta<\pi/2\)。结合贴面方案失败的条件，得到

\[\boxed{\arccos\frac23<\theta<\frac\pi2}. \tag{14}\]

若把两种方案的临界状态都计入备选，则可把左端写为闭区间；但在
\(\theta=\arccos(2/3)\)
时，速度更小的贴面方案已经可行，通常把切线抛射方案的左端写成开区间。

还需检查小球确实能在发射点脱离。若假想它仍沿圆弧运动，所需径向加速度为
\(v^2/R\)。由式（12）

\[\frac{v^2}{R}-g\cos\theta
=g\left[\frac{1+\cos\theta}{2\cos\theta}-\cos\theta\right]
=\frac{g(1-\cos\theta)(1+2\cos\theta)}{2\cos\theta}>0,\]

说明重力不足以提供所需向心力，曲面也不能施加向内拉力，小球会立即离面，方案自洽。

回看整个推导，四问实际上由两个互相独立的判据连接起来。能量守恒只负责给出``若小球能沿圆弧走到坡顶，初速度应当多大''；径向动力学则负责判断曲面能否提供实现这条轨迹所需的约束力。列径向方程时，应先在图上把``指向球心''规定为正方向。重力的径向分量指向球心，而外表面的支持力背离球心，因此有

\[mg\cos\varphi-N=\frac{mv^2}{R},\]

而不是把两项相加。把能量关系代入后，\(N\) 成为 \(\cos\varphi\)
的一次函数。上滑过程中 \(\varphi\) 逐渐减小、\(\cos\varphi\)
逐渐增大，所以支持力的最小值就在发射点；这说明只需检查初态，而不是凭感觉检查坡顶。临界值
\(N_0=0\) 既是贴面方案的右端边界，也是切线抛射方案开始出现的分界。

切线抛射部分要把几何和运动学逐项对应：初速度方向与水平面的夹角就是
\(\theta\)，故
\(v_x=-v\cos\theta\)、\(v_y=v\sin\theta\)；小洞相对发射点的水平位移和竖直位移分别为
\(R\sin\theta\) 与
\(R(1-\cos\theta)\)。由水平方程先消去飞行时间，再代入竖直方程，可以避免同时求
\(t\) 和
\(v\)。最后再用``假想仍贴面时需要多大的向心力''检查立即离面，才算把新方案的几何命中条件和动力学可实现性同时闭合。

代数上可令 \(c=\cos\theta\) 检查两分支的衔接。临界处
\(c=2/3\)，贴面方案给出 \(v^2=2gR(1-c)=2gR/3\)，切线方案给出
\(v^2=gR(1+c)/(2c)=5gR/4\)。两种速度并不相等，因为它们对应不同的初始运动状态；临界处贴面方案已经可行且更省速度，所以切线方案的左端应取开区间。这个比较表明：方案分界由约束可行性确定，两个优化目标未必在边界连续。

贴面与脱离两种方案必须用支持力条件分界：

\begin{enumerate}
\def\labelenumi{\arabic{enumi}.}
\tightlist
\item
  能量方程只给速率，径向方程才给支持力。对外球面取指向球心为正时，重力分量与支持力方向相反，符号必须按图
  1-2 写清。
\item
  把``到顶速度为零''的初速度代入 \(N(\varphi)\) 后，要在整个
  \(0\le\varphi\le\theta\) 上求最小值，而不是只检查坡顶。
\item
  脱离后重新建立抛体坐标，发射速度沿初始切线。分别写水平、竖直位移并消去飞行时间，才能得到经过坡顶的速度条件。
\item
  两种方案的临界角应在同一点衔接，并检查发射初瞬所需的支持力不为正；这同时验证区间端点的开闭。
\end{enumerate}

\textbf{支持力与脱离条件：} 贴面阶段应始终并列保留能量式与径向式：

\[\begin{aligned}
v^2(\varphi)&=v_0^2-2g\,\Delta h(\varphi),\\
mg\cos\varphi-N(\varphi)&=\frac{mv^2(\varphi)}{R},\\
N(\varphi)&=mg\cos\varphi-\frac{m}{R}
\left[v_0^2-2g\Delta h(\varphi)\right].
\end{aligned}\]

能量式只决定速度；只有第三行在整个运动区间内满足
\(N\ge0\)，才说明所假设的贴面运动确实能够实现。

\textbf{轨迹与山坡之间的间隙。} 图 1-2 右图以
\(R=g=1\)、\(\theta=60^\circ\)
为例，曲线由抛体方程直接绘出。命中顶点还须保证之前没有再撞山。令
\(c=\cos\theta\)、\(s=\sin\theta\)，以无量纲飞行参数
\(u=v\cos\theta\,t/R\) 表示轨迹：

\[x/R=s-u,\qquad y/R=c+\frac{s}{c}u-\frac{u^2}{c(1+c)}.\]

当 \(0<u<s\) 时，

\[\frac{x^2+y^2-R^2}{R^2}=\frac{u^2(u-s)[u-s(1+2c)]}{c^2(1+c)^2}>0.\]

两个括号都为负，因而球始终在球面外，直到 \(u=s\)
落入洞口。贴面方案``恰好到顶且速度为零''还存在理想模型的时间问题：近顶端
\(v\approx\sqrt{gR}\,\varphi\)，所需时间含
\(\sqrt{R/g}\int d\varphi/\varphi\)，发散。因此式（2）是能量临界速度；实际有限小洞或略大的初速度才可在有限时间入洞。

\subsection{六、关键洞察}\label{ux516dux5173ux952eux6d1eux5bdf}

能量方程只决定``若沿这条轨迹运动，速度应是多少''，不能保证约束力有物理意义。外圆弧问题必须把
\(N\geq0\)
与能量条件并列检查。后两问的本质是：当约束轨迹不可实现时，利用同一点的切线方向构造一条无约束抛物线。

\subsection{七、易错点}\label{ux4e03ux6613ux9519ux70b9}

\begin{itemize}
\tightlist
\item
  把径向方程写成 \(N+mg\cos\varphi=mv^2/R\)。对外圆弧而言，\(N\)
  指向外侧，应带负号。
\item
  只在坡顶检查支持力。支持力最小处是初始点，不是坡顶。
\item
  得到式（2）后直接认定方案可行，遗漏 \(N\geq0\)。
\item
  在切线抛射中把水平速度写成 \(v\sin\theta\)；切线与水平方向的夹角是
  \(\theta\)，其水平分量大小为 \(v\cos\theta\)。
\item
  把切线抛射也误设成``到顶速度为零''。抛体击中坡顶时通常仍有速度。
\end{itemize}

\subsection{八、课堂讲解建议}\label{ux516bux8bfeux5802ux8bb2ux89e3ux5efaux8bae}

先只让学生用能量求式（2），随后追问：``半球表面能不能拉住小球？''再引入径向方程。这样可以清楚区分运动学上写出的圆弧轨迹与动力学上真正可实现的轨迹。讲抛射方案时，应让学生先在图上标出切线单位向量，再写两个分方向方程。

板书时建议把页面分成``能量条件''和``接触条件''两列。左列写出速度随位置的关系，右列固定向心正方向并列径向方程。两列合并以后再讨论
\(N\) 的最小位置，学生会清楚看到：能量是必要条件，\(N\geqslant0\)
才决定轨迹能否实现。可用``桌面能推物体但不能拉物体''帮助建立支持力不能为负的直觉。

进入第（3）问时不要直接公布切线抛射。先让学生判断贴面方案失败后小球会向半球内还是半球外运动，再在发射点画出切线。要求学生逐项写出小洞相对发射点的横、纵坐标，完成消元后用量纲和
\(\theta\to0\)
的极限检查公式。最后再回到径向方程证明立即离面，使课堂形成``提出新轨迹---算出命中条件---验证动力学自洽''的完整闭环。

\begin{itemize}
\tightlist
\item
  基础层学生重点掌握能量方程与支持力符号；
\item
  进阶层学生应能说明为什么支持力最小值出现在发射点；
\item
  提高层学生可比较临界角处两种方案的速度，并解释为何优先采用贴面方案。
\end{itemize}

\subsection{九、变式训练}\label{ux4e5dux53d8ux5f0fux8badux7ec3}

\begin{enumerate}
\def\labelenumi{\arabic{enumi}.}
\tightlist
\item
  若坡顶小洞位于圆心角 \(\beta\)
  处而不是顶端，推导切线抛射命中该洞的初速度。
\item
  若半球表面存在足够大的静摩擦，讨论能否扩大贴面方案的角度范围，并说明摩擦力对径向方程是否有贡献。
\end{enumerate}

第一题训练把几何位移从特定顶点推广到任意目标点。解题时先用两个半径方向的坐标差写出目标位移，再用切线速度分解消去飞行时间；最后检查所得时间为正。第二题训练学生区分切向力和径向力：静摩擦能改变速度---位置关系，却不直接出现在径向方程中，因此必须重新联立能量或功---能关系与
\(N\geqslant0\)，不能只在原结论上机械修改。

还可把半球改为半径随位置缓慢变化的凸曲面，让学生用曲率半径 \(\rho\) 代替
\(R\)
写临界离面条件。该变式的目的不是完成复杂计算，而是辨认``法向约束失效''这一模型可以迁移到任意光滑凸轨道。

\textbf{变式解答。} 洞口位于 \(\beta<\theta\)。水平位移大小
\(D=R(\sin\theta-\sin\beta)\)，高度差
\(Y=R(\cos\beta-\cos\theta)\)。沿原切线发射时

\[t=\frac{D}{v\cos\theta},\qquad v^2=\frac{gD^2}{2\cos^2\theta(D\tan\theta-Y)}.\]

要求
\(D>0\)、分母为正，再检查中间轨迹在球面外。对质点光滑坡面，静摩擦不能任意调节能量；若改为有转动惯量的滚动球，须另给纯滚动条件及转动动能，不能仅在原质点结论中加入一个摩擦力。

\newpage

\section{第 2 题}\label{ux7b2c-2-ux9898}

在坐标原点以初速度 \(v_0\) 斜向上抛出一质点，初速度与水平方向夹角为
\(\theta\)，忽略空气阻力，重力加速度为 \(g\)。

\begin{enumerate}
\def\labelenumi{\arabic{enumi}.}
\tightlist
\item
  质点需要到达点 \((l,0)\)。两点之间有一座山，山体上边界为
  \[y_h(x)=y_0-\frac{k}{2}\left(x-\frac l2\right)^2,\] 且
  \(y_0-kl^2/8<0\)。要求抛物线不与山体相交，求最小初速度，并给出参数之间相应的分类条件。
\item
  另一座山的上边界为 \[y_h(x)=y_0-\frac{1}{2l}(x-x_0)^2,\] 且
  \(y_h(0)<0\)、\(y_h(l)<0\)。要求向右抛出的质点在飞行过程中不与山体相交。求最小初速度及相应的分类条件；以
  \(x_0\) 为横轴、\(y_0\) 为纵轴，标出不同情形所在的参数区域。
\end{enumerate}

\subsection{一、题型定位}\label{ux4e00ux9898ux578bux5b9aux4f4d-1}

在避开抛物线山体的轨迹中求最小初速。固定落点时只需优化发射角；落点自由时，则按山体高度判断最优轨迹经过山脚还是与山体相切。

\subsection{二、知识点}\label{ux4e8cux77e5ux8bc6ux70b9-1}

\begin{itemize}
\tightlist
\item
  斜抛轨迹方程及参数化 \(\lambda=\tan\theta\)。
\item
  固定射程条件下 \(v_0^2=gl(1+\lambda^2)/(2\lambda)\)。
\item
  两条抛物线的高低比较和二次函数最小值。
\item
  临界不相交对应相切，即二次方程出现重根。
\item
  有约束极值、参数区域分类和物理解的筛选。
\end{itemize}

\subsection{三、物理图像}\label{ux4e09ux7269ux7406ux56feux50cf-1}

固定起点和落点后，所有抛物线都可写成 \(y=\lambda x(1-x/l)\)。当
\(\lambda\) 增大时，轨迹整体抬高，但所需速度在 \(\lambda=1\)
前先减小、之后增大，所以``能越过山的最小
\(\lambda\)''并不总等于``最小速度''；只有约束迫使 \(\lambda>1\)
时才取临界轨迹。

第二问中，若山体较宽缓，用最省射程速度的 \(45^\circ\)
轨迹落在右山脚即可；若山体较陡窄，这条轨迹会穿山，最优轨迹必在山坡某点相切。若最优轨迹与山体之间仍有正间隙，就还能略微降低速度，因此不可能是极小值。

\subsection{四、解题入口}\label{ux56dbux89e3ux9898ux5165ux53e3-1}

统一把斜率参数记为

\[\lambda=\tan\theta>0.\]

一般斜抛轨迹为

\[y_p(x)=\lambda x-\frac{g(1+\lambda^2)}{2v_0^2}x^2. \tag{1}\]

第一问先用 \(y_p(l)=0\) 消去
\(v_0\)；第二问则保留二次项系数，并用``重根''表示相切。

\subsection{五、解析}\label{ux4e94ux89e3ux6790-1}

图 2-1 左侧比较固定落点 \((l,0)\)
时的抛物线与抛物型山体，最小初速对应二次函数间的最小安全间隙；右侧表示落点不固定时的两类临界轨迹：恰好受山脚约束，或与山坡相切。第（1）步使用左图，第（2）步根据右图分类求最优条件。

\begin{figure}[H]
\centering
\includegraphics[width=0.68\linewidth,height=\textheight,keepaspectratio,alt={图 2-1　固定落点时轨迹与山体的比较，以及第二问的两类最优轨迹。}]{../figures/solution-02.pdf}
\caption*{图 2-1　固定落点时轨迹与山体的比较，以及第二问的两类最优轨迹。}
\end{figure}

\textbf{（1）固定到达 \((l,0)\)}

按图 2-1 左侧，轨迹与山体都是二次函数。先用固定落点消去
\(v_0\)，再要求整个区间内``轨迹高度减山体高度''非负；临界最优状态就是这个二次差函数的最小值恰为零。

令轨迹经过 \((l,0)\)。把 \(x=l\) 代入式（1）：

\[0=\lambda l-\frac{g(1+\lambda^2)}{2v_0^2}l^2.\]

因此

\[v_0^2=\frac{gl(1+\lambda^2)}{2\lambda}
=\frac{gl}{\sin2\theta}. \tag{2}\]

无约束最小值也可直接在 \(\lambda\) 变量中证明。把式（2）写成

\[v_0^2=\frac{gl}{2}\left(\lambda+\frac1\lambda\right).\]

因向右上方发射有 \(\lambda>0\)，所以

\[\lambda+\frac1\lambda\geq2,\]

等号仅在

\[\lambda=1\]

时成立。于是

\[v_{0,\mathrm{free}}^2=gl,\]

并由 \(\lambda=\tan\theta\) 得
\(\theta=45^\circ\)。这一证明同时说明自由极值是唯一的全局最小值，而不是只由``最大射程角为
\(45^\circ\)''类比得到。

同时，轨迹可写成

\[y_p(x)=\lambda x\left(1-\frac{x}{l}\right). \tag{3}\]

令 \(z=x-l/2\)。利用

\[x\left(1-\frac{x}{l}\right)=\frac l4-\frac{z^2}{l},\]

可得

\[y_p=\frac{\lambda l}{4}-\frac{\lambda}{l}z^2,\]

而山体为

\[y_h=y_0-\frac{k}{2}z^2.\]

两者之差为

\[D(z)=y_p-y_h
=\left(\frac{\lambda l}{4}-y_0\right)
+\left(\frac k2-\frac\lambda l\right)z^2. \tag{4}\]

题设 \(y_0-kl^2/8<0\) 保证山脚在两个端点之间，且在 \(x=0,l\)
处山体低于地面。

若没有山体约束，式（2）在 \(\lambda=1\)，即 \(\theta=45^\circ\)
时最小，得到

\[v_{0,\mathrm{free}}=\sqrt{gl}. \tag{5}\]

\(45^\circ\) 轨迹的最高点在 \(x=l/2\)，高度为 \(l/4\)。

若

\[y_0\leq\frac l4, \tag{6}\]

则式（4）在整个区间均非负：当二次项系数非负时最小值在
\(z=0\)，由式（6）保证不撞山；当二次项系数为负时最小值在端点，题设已保证端点安全。因此自由极值可实现：

把 \(\lambda=1\) 代入差函数，得到

\[D(z)=\left(\frac l4-y_0\right)
+\left(\frac k2-\frac1l\right)z^2.\]

若 \(k/2-1/l\geq0\)，则 \(D(z)\) 在 \(z=0\) 最小：

\[D_{\min}=D(0)=\frac l4-y_0\geq0.\]

若 \(k/2-1/l<0\)，则 \(D(z)\) 随 \(|z|\)
增大而减小，在实际比较区间的端点处最小；题设保证山体在两端低于地面，而抛物线两端位于地面，故端点差仍非负。两种开口关系都已覆盖，才可以确认整段轨迹不穿山。

\[\boxed{v_{0\min}=\sqrt{gl},\qquad\theta=45^\circ}. \tag{7}\]

若

\[y_0>\frac l4, \tag{8}\]

则 \(45^\circ\) 轨迹在山顶下方。要越过山顶，必须满足

\[\frac{\lambda l}{4}\geq y_0.\]

因为此时临界值 \(\lambda=4y_0/l>1\)，而式（2）在 \(\lambda>1\) 上随
\(\lambda\) 增大而增大，所以最小速度对应恰好擦过山顶：

\[\lambda_*=\frac{4y_0}{l}. \tag{9}\]

题设 \(y_0<kl^2/8\) 等价于
\(\lambda_*<kl/2\)，保证式（4）的最小点确实在山顶。代入式（2）：

\[\boxed{
v_{0\min}
=\sqrt{\frac{gl}{\sin\left(2\arctan\frac{4y_0}{l}\right)}}
}. \tag{10}\]

相应发射角为

\[\boxed{\theta_*=\arctan\frac{4y_0}{l}}. \tag{11}\]

\textbf{（2）落点不固定时的山体分类}

图 2-1
右侧说明最优轨迹的约束可能在端点激活，也可能在山坡内部激活。前者让抛物线恰过山脚，后者让轨迹与山体相切；相切条件应写成位置相同且斜率相同，等价于二次差函数具有重根。

先讨论有实际山峰的情形
\(y_0>0\)，并设山位于发射点右侧。山体与地面的两个交点为

\[x_\pm=x_0\pm\sqrt{2ly_0}.\]

记右山脚

\[x_R=x_0+\sqrt{2ly_0}. \tag{12}\]

第一类情形是用 \(45^\circ\) 轨迹恰好落在右山脚。对给定射程
\(x_R\)，最小速度为

\[v_0=\sqrt{gx_R},\]

轨迹为

\[y_{45}(x)=x-\frac{x^2}{x_R}. \tag{13}\]

记半个山脚宽度 \(b=\sqrt{2ly_0}\)，则
\(x_R=x_0+b\)。轨迹和山体在右山脚相交，差函数可准确分解为

\[y_{45}-y_h=(x_R-x)\left[\frac{x_0-b}{2l}+\left(\frac1{x_R}-\frac1{2l}\right)x\right].\]

山位于右侧时 \(x_0>b\)，括号是关于 \(x\) 的一次函数。它在 \(x=0\)
为正，在 \(x=x_R\) 等于 \(1-b/l\)。所以整段不穿山的充要条件为

\[b\le l\quad\Longleftrightarrow\quad y_0\le\frac l2. \tag{14}\]

这不是 \(x_R<2l\)；后者只是一个较强的充分条件。于是第一分支为

\[\boxed{v_{0\min}=\sqrt{g(x_0+\sqrt{2ly_0})},\qquad\theta=45^\circ,\quad0<y_0\le l/2.} \tag{15}\]

第二分支为 \(y_0>l/2\)。此时在右山脚附近 \(y_{45}-y_h<0\)，

此时式（13）不能保证越过山体，最优轨迹应与山体相切。把一般轨迹式（1）与山体联立：

\[\left(\frac1{2l}-\frac{g(1+\lambda^2)}{2v_0^2}\right)x^2
+\left(\lambda-\frac{x_0}{l}\right)x
+\left(\frac{x_0^2}{2l}-y_0\right)=0. \tag{17}\]

相切意味着式（17）有重根，判别式为零。令

\[D_0=x_0^2-2ly_0>0,\]

化简得到

\[v_0^2(\lambda)
=\frac{gD_0(1+\lambda^2)}
{2x_0\lambda-l\lambda^2-2y_0}. \tag{18}\]

式（18）的来龙去脉值得展开。把式（17）记成

\[Ax^2+Bx+C=0,\]

其中

\[A=\frac1{2l}-\frac{g(1+\lambda^2)}{2v_0^2},\]

\[B=\lambda-\frac{x_0}{l},\]

\[C=\frac{x_0^2-2ly_0}{2l}=\frac{D_0}{2l}.\]

相切条件 \(B^2-4AC=0\) 给出

\[\left(\lambda-\frac{x_0}{l}\right)^2
-\frac{2D_0}{l}\left[
\frac1{2l}-\frac{g(1+\lambda^2)}{2v_0^2}
\right]=0.\]

把含 \(v_0\) 的项留在一侧：

\[\frac{gD_0(1+\lambda^2)}{lv_0^2}
=\frac{D_0}{l^2}-\left(\lambda-\frac{x_0}{l}\right)^2.\]

右端展开为

\[\begin{aligned}
\frac{D_0}{l^2}-\left(\lambda-\frac{x_0}{l}\right)^2
&=\frac{x_0^2-2ly_0}{l^2}
-\lambda^2+\frac{2x_0\lambda}{l}-\frac{x_0^2}{l^2}\\
&=\frac{2x_0\lambda-l\lambda^2-2y_0}{l}.
\end{aligned}\]

交叉相乘即得到式（18）。因此分母为正并非附加的形式要求，而是 \(v_0^2>0\)
的物理必要条件。

分母必须为正，这同时给出允许的 \(\lambda\) 区间。对式（18）求极值。令
\(v_0^2\) 的导数为零，分子中的公共正因子约去后得到

\[x_0\lambda^2+(l-2y_0)\lambda-x_0=0. \tag{19}\]

下面把求导步骤完整写出。设

\[P(\lambda)=1+\lambda^2,\]

\[Q(\lambda)=2x_0\lambda-l\lambda^2-2y_0.\]

则

\[v_0^2=gD_0\frac{P}{Q},\]

且极值条件为

\[P'Q-PQ'=0.\]

代入 \(P'=2\lambda\)、\(Q'=2x_0-2l\lambda\)：

\[2\lambda(2x_0\lambda-l\lambda^2-2y_0)
-(1+\lambda^2)(2x_0-2l\lambda)=0.\]

两边同除以 \(2\)，逐项展开并合并同类项：

\[\begin{aligned}
0={}&\lambda(2x_0\lambda-l\lambda^2-2y_0)
-(1+\lambda^2)(x_0-l\lambda)\\
={}&x_0\lambda^2+(l-2y_0)\lambda-x_0.
\end{aligned}\]

三次项自动相消，说明把 \(\lambda=\tan\theta\)
作为优化变量确实能显著降低代数复杂度。

在 \(x_0>0\) 的物理区域内取正根：

\[\boxed{
\lambda_*
=\frac{2y_0-l+\sqrt{(l-2y_0)^2+4x_0^2}}{2x_0}
}. \tag{20}\]

于是

\[\boxed{
v_{0\min}
=\sqrt{
\frac{g(x_0^2-2ly_0)(1+\lambda_*^2)}
{2x_0\lambda_*-l\lambda_*^2-2y_0}
}
}. \tag{21}\]

相切点由式（17）的重根给出。若需要核查轨迹与山体接触的位置，可用

\[x_t=-\frac{\lambda_*-x_0/l}
{2\left[1/(2l)-g(1+\lambda_*^2)/(2v_{0\min}^2)\right]}. \tag{22}\]

应选 \(x_t>0\) 且位于山体地面以上区间内的解。

为了在 \((x_0,y_0)\) 平面上作图，令

\[X=\frac{x_0}{l},\qquad Y=\frac{y_0}{l}.\]

题设条件写成

\[Y<\frac{X^2}{2},\qquad Y<\frac{(1-X)^2}{2}. \tag{23}\]

右侧有山时还要求 \(X>\sqrt{2Y}\)、\(Y>0\)。两分支的分界是

\[\boxed{Y=\frac12}. \tag{24}\]

在题设有效区域 \(Y<X^2/2\)、\(Y<(1-X)^2/2\) 内，\(0<Y\le1/2\)
用右山脚分支，\(Y>1/2\) 用相切分支。图 2-1
下图把这条水平分界和两条题设边界同时标出；边界曲线之外的区域不满足原题条件。若
\(y_0\leq0\)，山体最高点不高于地面；若整座山位于发射点左侧，则向右抛射不会遇到山体，这两类情形的初速度下确界为零，通常不属于题目希望讨论的越山情形。

这类最优化题不能一开始就假设``最优轨迹必与山相切''。固定落点时，先由落点条件把
\(v_0^2\) 写成 \(\lambda=\tan\theta\) 的函数，再对 \(\lambda\)
求无约束最小值；只有当这条最省速度的 \(45^\circ\)
轨迹穿山时，山体约束才会起作用。比较抛体轨迹 \(y_p(x)\) 与山体
\(y_h(x)\) 时，应研究差函数

\[F(x)=y_p(x)-y_h(x).\]

要求不相交就是在山体所在区间内
\(F(x)\geqslant0\)。临界状态的最小值恰好降到零，因此同时满足
\(F(x_*)=0\)、\(F'(x_*)=0\)；这就是相切条件的严格来源。若差函数为二次式，也可等价地令判别式为零，但仍须检查重根
\(x_*\) 是否位于真实山坡区间内。

第二问没有固定终点，因此还多出一种边界机制：轨迹可以在右山脚之后落地。若到右山脚的
\(45^\circ\)
抛射已经越过整座山，它就是无约束最优解；否则必须提升轨迹直到与山坡相切。参数图中的分界线
\(y_0=l/2\) 对应 \(45^\circ\)
轨迹在右山脚的斜率刚好与山坡相同。完成代数后，应依次检查分母为正、\(\lambda>0\)、相切点在山体上以及题设给出的两个山脚条件，不能只保留二次方程的正根。

数学上建议先无量纲化：\(X=x_0/l\)、\(Y=y_0/l\)、\(u=v_0/\sqrt{gl}\)。这样分类边界只在
\((X,Y)\) 平面中出现，不再混杂长度单位。相切分支中消元得到的关于
\(\lambda\) 的二次方程有两个代数根，可用 \(\lambda>0\)
和轨迹向右上方发射筛去非物理解。把临界根代回差函数后，\(F(x)\)
应可写成正系数乘
\((x-x_*)^2\)；这比只检验判别式为零更直接地证明整条轨迹不低于山体。

越山最优化要把``经过目标''和``全程不穿山''分开：

\begin{enumerate}
\def\labelenumi{\arabic{enumi}.}
\tightlist
\item
  先从 \(x=v_0\cos\theta\,t\) 消去
  \(t\)，逐步得到轨迹方程；固定落点时再用 \(y(l)=0\) 约束发射角。
\item
  山体约束要求轨迹减山体在相关区间内非负。宽缓情形的最优轨迹由
  \(45^\circ\) 最省射程给出，但仍要检查山顶高度。
\item
  相切情形把两条二次曲线相减得到关于 \(x\) 的二次式。完整写出 \(A,B,C\)
  后令 \(B^2-4AC=0\)，再解出 \(v_0^2(\lambda)\)。
\item
  对 \(\lambda\)
  求极值时应写分子、分母导数和允许分母为正的范围；求得驻点后还要检查切点位于山体实际区间，而不是代数延长线上。
\end{enumerate}

\textbf{轨迹间距：} 抛物线轨迹应先由时间参数消元：

\[\begin{aligned}
x&=v_0\cos\theta\,t,
&t&=\frac{x}{v_0\cos\theta},\\
y(x)&=x\tan\theta-
\frac{gx^2}{2v_0^2\cos^2\theta}.
\end{aligned}\]

若山体轮廓为二次函数 \(y_m(x)\)，令
\(F(x)=y(x)-y_m(x)=Ax^2+Bx+C\)，临界不穿山条件就是

\[F(x_*)=0,\qquad F'(x_*)=0
\quad\Longleftrightarrow\quad B^2-4AC=0.\]

\textbf{分界的反例检查。} 取 \(x_0=3l\)、\(y_0=l/8\)，有
\(b=l/2\)、\(x_R=7l/2>2l\)。原有按 \(x_R\)
的分类会误判为相切分支；准确差函数的两个端点均为正，\(45^\circ\)
轨迹完全越山。取 \(y_0=l/2\) 时，式（19）给 \(\lambda_*=1\)，式（21）给
\(v_0^2=g(x_0+l)\)，与右山脚分支连续衔接。若严格禁止任何接触，切触或落在山脚的速度应理解为下确界，实际发射速度需略大。

\subsection{六、关键洞察}\label{ux516dux5173ux952eux6d1eux5bdf-1}

第一问的优化顺序应是``先判断无约束极值是否可行，再寻找刚好满足约束的临界轨迹''。第二问的两个答案分支对应两种不同的临界机制：落点受右山脚约束，或轨迹受山坡切触约束。相切并不是经验猜测，而是由``若有余量还能继续降速''严格推出。

\subsection{七、易错点}\label{ux4e03ux6613ux9519ux70b9-1}

\begin{itemize}
\tightlist
\item
  把斜抛轨迹的二次项写成 \(gx^2/(2v_0^2\cos\theta)\)，漏掉一个
  \(\cos\theta\)。
\item
  第一问只比较最高点而不检查其他位置。必须借助式（4）的二次函数结构说明全区间安全。
\item
  在 \(y_0>l/4\) 时仍取 \(45^\circ\)，忽略山顶高度约束。
\item
  第二问把``相交''误当成临界；最优轨迹与山坡应为重根相切。
\item
  求出 \(\lambda_*\) 后不检查式（18）的分母、相切点位置和发射方向。
\item
  混淆参数平面的纵坐标 \(y_0\) 与实际轨迹纵坐标 \(y\)。
\end{itemize}

\subsection{八、课堂讲解建议}\label{ux516bux8bfeux5802ux8bb2ux89e3ux5efaux8bae-1}

先画三条固定端点抛物线，让学生观察 \(\lambda\)
改变时高度和速度并不同步。第二问可先让学生解释为什么最优轨迹必``碰到一个约束边界''，再分别指出右山脚和山坡切点。分类图中应把题设允许区域与方案分界线同时画出。

第一问建议先暂时遮住山体，只解决``从原点到 \((l,0)\)
的最小发射速度''。得到 \(45^\circ\)
后再恢复山体，检查这条无约束最优轨迹是否安全。这样学生能够理解分类讨论不是人为拆分，而是无约束极值是否满足不等式约束的自然结果。比较两条抛物线时，应始终把
\(F(x)=y_p-y_h\) 写在板书中央，并用 \(F\geqslant0\) 表示``不穿山''。

第二问先在参数图上画出两个山脚，再讨论最优轨迹可能触及哪一个边界。对于相切分支，要求学生写出
\(F(x_*)=F'(x_*)=0\)，并解释它与判别式为零的等价性。求出 \(\lambda\)
后安排一次``物理解筛选''：分母为正、发射角向右、相切点位于山坡、落点在右山脚之外，四项缺一不可。

课堂末尾可让学生用一句话概括两种分界机制：``先到右山脚''和``先碰山坡''。若学生只能记住复杂根式而说不出这两个图像，说明分类逻辑还没有建立。

\subsection{九、变式训练}\label{ux4e5dux53d8ux5f0fux8badux7ec3-1}

\begin{enumerate}
\def\labelenumi{\arabic{enumi}.}
\tightlist
\item
  将第一问山顶从 \(x=l/2\) 平移到
  \(x=a\)，讨论固定落点时最优轨迹是否仍在山顶相切。
\item
  在第二问中限定发射角 \(0<\theta\leq\theta_m\)，重新判断 \(45^\circ\)
  分支和相切分支的最优解。
\end{enumerate}

第一题训练非对称约束下的活跃点判断。应先写一般差函数，再由 \(F=F'=0\)
求真实切点；山顶不再位于中点时，最危险位置通常不会凭对称性直接确定。第二题把角度上限加入优化问题：若无约束最优角或相切角超过
\(\theta_m\)，新的最优解会落在角度边界 \(\theta=\theta_m\) 上。

还可在给定初速度后反求最大山高，训练把同一约束改写为允许参数上限。

\textbf{变式解答。} 在固定射程问中取 \(y_0=l/2\)、\(k=5/l\)，满足
\(y_0-kl^2/8=-l/8<0\)。临界斜率 \(\lambda_*=2\)，故
\(v_{0\min}^2=5gl/4\)，\(\theta_* =\arctan2\)。差函数为
\((1/2l)(x-l/2)^2\)，除了山顶切点外均为正，直接验证全程安全。

\newpage

\section{第 3 题}\label{ux7b2c-3-ux9898}

一单摆由长度为 \(l\)、不带电且不可伸长的细绳和质量为 \(m\)、电量为 \(q\)
的带电质点组成。摆的下方有一接地金属平板。真空静电力常量为
\(k\)，重力加速度为 \(g\)。初始时摆线与竖直方向夹角为
\(\alpha_0=\pi/3\)，质点到金属板的垂直距离为 \(d=3l/4\)。

\begin{enumerate}
\def\labelenumi{\arabic{enumi}.}
\tightlist
\item
  求初始时质点所受合外力的大小和方向。
\item
  求初始时绳子的张力，以及质点第一次到达最低点时绳子的张力。
\item
  质点最终在最低点附近做小振动，求小振动周期 \(f_3\)。
\item
  若把金属平板换成半径为 \(R\) 的接地金属球，其中
  \(R>l\)，且金属球顶点位于悬点在原金属板上的垂足处，求此时最低点附近的小振动周期
  \(f_4\)。
\end{enumerate}

\begin{figure}[H]
\centering
\includegraphics[width=0.761\linewidth,height=\textheight,keepaspectratio,alt={图 3-1　带电单摆位于接地金属板上方。}]{../figures/question-03.pdf}
\caption*{图 3-1　带电单摆位于接地金属板上方。}
\end{figure}

\subsection{一、题型定位}\label{ux4e00ux9898ux578bux5b9aux4f4d-2}

用接地导体的镜像电荷求电力，再处理摆球运动。释放速度由势能差确定，小振动周期由最低点附近的线性恢复力确定。

\subsection{二、知识点}\label{ux4e8cux77e5ux8bc6ux70b9-2}

\begin{itemize}
\tightlist
\item
  接地无限平板的镜像电荷及静电势能 \(U_e=-kq^2/(4h)\)。
\item
  接地球的像电荷法及能量中的 \(1/2\) 因子。
\item
  单摆的径向、切向动力学和机械能守恒。
\item
  小振动的势能二阶展开。
\item
  广义坐标 \(\alpha\) 下的动能 \(T=\frac12ml^2\dot\alpha^2\)。
\end{itemize}

\subsection{三、物理图像}\label{ux4e09ux7269ux7406ux56feux50cf-2}

接地平板对带电摆球的作用等效为平板另一侧、与摆球等距的异号像电荷。像电荷吸引力始终垂直平板向下，所以与重力方向相同，但其大小随摆球到板距离按
\(h^{-2}\) 变化。

在最低点附近，摆球的水平位移是一阶小量，而它到导体的距离变化是二阶小量。因此恢复力的线性项要从势能对摆角的二阶项中读取，平板情形在线性近似中可用最低点的静电力修正重力；接地球还会产生随横向位移变化的水平力，须另算势能曲率。

\subsection{四、解题入口}\label{ux56dbux89e3ux9898ux5165ux53e3-2}

先由初始几何求悬点到平板的高度

\[H=d+l\cos\alpha_0=\frac54l.\]

摆角为 \(\alpha\) 时，摆球到平板的距离为

\[h(\alpha)=H-l\cos\alpha.\]

平板情形统一使用

\[F_e(h)=\frac{kq^2}{4h^2},\qquad U_e(h)=-\frac{kq^2}{4h}.\]

接地球情形则先写球心到摆球的距离 \(s(\alpha)\)，再写镜像势能。

\subsection{五、解析}\label{ux4e94ux89e3ux6790-2}

图 3-2 左侧用像电荷 \(-q\)
表示接地平板的静电作用；中图画出摆球在最低点所受张力、重力与静电吸引力；右侧给出接地球情形中悬点到球心距离
\(a=R+5l/4\)、摆长 \(l\) 与摆球到球心距离 \(s\)
的几何关系。第（1）---（3）步使用前两图，第（4）步使用右图建立镜像势能。

\begin{figure}[H]
\centering
\includegraphics[width=0.9\linewidth,height=\textheight,keepaspectratio,alt={图 3-2　平板镜像电荷、最低点受力及接地球情形的几何量。}]{../figures/solution-03.pdf}
\caption*{图 3-2　平板镜像电荷、最低点受力及接地球情形的几何量。}
\end{figure}

\textbf{（1）初始合外力}

按图 3-2 左侧，真实电荷到接地平板的距离为
\(h_0\)，像电荷在板下对称位置，故两电荷间距是
\(2h_0\)。先求竖直向下的静电力，再与重力合成并投影到摆线切向。

初始时摆球到平板的距离为

\[h_0=d=\frac34l.\]

像电荷与摆球相距 \(2h_0\)，所以静电吸引力大小为

\[F_{e0}=\frac{kq^2}{(2h_0)^2}
=\frac{kq^2}{4h_0^2}
=\frac{4kq^2}{9l^2}. \tag{1}\]

重力和静电力都竖直向下，合成竖直力大小

\[F_0=mg+\frac{4kq^2}{9l^2}. \tag{2}\]

初始时由静止释放，径向速度为零，因而径向加速度 \(v^2/l=0\)。绳张力抵消
\(F_0\) 的径向分量，合外力只剩切向分量：

\[F_{\text{合},0}=F_0\sin\alpha_0
=\frac{\sqrt3}{2}\left(mg+\frac{4kq^2}{9l^2}\right). \tag{3}\]

方向沿圆弧切线指向最低点。因此

\[\boxed{\boldsymbol F_{\text{合},0}\text{ 沿切线指向最低点，大小为式（3）}}.\]

\textbf{（2）初始张力与最低点张力}

图 3-2
中图显示最低点的重力与静电力同向向下，张力向上。最低点张力除平衡两种向下力外，还必须提供向心力
\(mv_b^2/l\)；速度 \(v_b\) 则由初态到最低点的能量守恒求出。

初始径向平衡给出

\[T_0=F_0\cos\alpha_0
=\frac12\left(mg+\frac{4kq^2}{9l^2}\right). \tag{4}\]

摆球到最低点时

\[h_b=H-l=\frac l4. \tag{5}\]

平板镜像体系的静电势能为

\[U_e(h)=-\frac{kq^2}{4h}. \tag{6}\]

从初始位置到最低点，重力势能减少

\[\Delta U_g=mgl(1-\cos\alpha_0)=\frac12mgl.\]

静电势能减少量为

\[\begin{aligned}
U_e(h_0)-U_e(h_b)
&=-\frac{kq^2}{3l}-\left(-\frac{kq^2}{l}\right)\\
&=\frac{2kq^2}{3l}.
\end{aligned}\]

能量守恒：

\[\frac12mv_b^2=\frac12mgl+\frac{2kq^2}{3l}. \tag{7}\]

所以

\[\frac{mv_b^2}{l}=mg+\frac{4kq^2}{3l^2}. \tag{8}\]

最低点处静电力大小为

\[F_{eb}=\frac{kq^2}{4(l/4)^2}=\frac{4kq^2}{l^2}. \tag{9}\]

沿竖直向上、即指向悬点的方向列径向方程：

\[T_b-mg-F_{eb}=\frac{mv_b^2}{l}.\]

代入式（8）、（9）：

\[\begin{aligned}
T_b
&=mg+\frac{4kq^2}{l^2}
+mg+\frac{4kq^2}{3l^2}\\
&=2mg+\frac{16kq^2}{3l^2}.
\end{aligned}\]

故

\[\boxed{T_0=\frac12\left(mg+\frac{4kq^2}{9l^2}\right)},\]

\[\boxed{T_b=2\left(mg+\frac{8kq^2}{3l^2}\right)}. \tag{10}\]

\textbf{（3）接地平板上方的小振动周期}

用最低点附近的小摆角 \(\alpha\) 作广义坐标。动能为

\[K=\frac12ml^2\dot\alpha^2. \tag{11}\]

摆球到平板的距离

\[\begin{aligned}
h(\alpha)
&=\frac54l-l\cos\alpha\\
&\approx\frac54l-l\left(1-\frac{\alpha^2}{2}\right)\\
&=\frac l4+\frac l2\alpha^2
=\frac l4(1+2\alpha^2).
\end{aligned} \tag{12}\]

重力势能相对最低点为

\[U_g=mgl(1-\cos\alpha)\approx\frac12mgl\alpha^2. \tag{13}\]

静电势能展开为

\[\begin{aligned}
U_e
&=-\frac{kq^2}{4h(\alpha)}\\
&=-\frac{kq^2}{l}\frac1{1+2\alpha^2}\\
&\approx-\frac{kq^2}{l}+\frac{2kq^2}{l}\alpha^2.
\end{aligned} \tag{14}\]

舍去与运动无关的常量，总势能二次项为

\[U(\alpha)=\frac12\left(mgl+\frac{4kq^2}{l}\right)\alpha^2. \tag{15}\]

与标准简谐振子

\[K=\frac12M_{\mathrm{eff}}\dot\alpha^2,\qquad
U=\frac12K_{\mathrm{eff}}\alpha^2\]

比较，得到

\[M_{\mathrm{eff}}=ml^2,\qquad
K_{\mathrm{eff}}=mgl+\frac{4kq^2}{l}.\]

因此

\[\boxed{
f_3=2\pi\sqrt{\frac{ml^2}{mgl+4kq^2/l}}
}. \tag{16}\]

题目把周期记作 \(f_3\)；若按通常记号用 \(T_3\) 表示周期，则右端相同。

\textbf{（4）接地金属球上方的小振动周期}

转看图 3-2 右侧。金属球的镜像电荷位置和大小都随摆球到球心的距离
\(s(\alpha)\) 改变，所以应先由余弦定理求
\(s^2=a^2+l^2-2al\cos\alpha\)，再把接地球的镜像势能对小角度展开。

金属球顶点位于原平板位置，故悬点到球心的距离为

\[a=R+\frac54l. \tag{17}\]

摆角为 \(\alpha\) 时，摆球到球心的距离满足余弦定理：

\[s^2(\alpha)=a^2+l^2-2al\cos\alpha. \tag{18}\]

最低点处

\[s_0=a-l=R+\frac l4. \tag{19}\]

对球外电荷 \(q\)，接地球的像电荷为

\[q'=-\frac{R}{s}q,\]

像电荷距球心 \(R^2/s\)。接地体系的静电能为实际电荷乘像电荷电势的一半：

\[U_e(s)=-\frac{kq^2R}{2(s^2-R^2)}. \tag{20}\]

在 \(\alpha=0\) 附近，\(\cos\alpha\approx1-\alpha^2/2\)，所以

\[s^2(\alpha)=s_0^2+al\alpha^2+O(\alpha^4). \tag{21}\]

记

\[D_0=s_0^2-R^2
=\left(R+\frac l4\right)^2-R^2
=\frac{l^2}{16}+\frac{lR}{2}. \tag{22}\]

则

\[\begin{aligned}
U_e(\alpha)
&=-\frac{kq^2R}{2(D_0+al\alpha^2)}\\
&\approx-\frac{kq^2R}{2D_0}
+\frac{kq^2Ral}{2D_0^2}\alpha^2.
\end{aligned} \tag{23}\]

加入重力势能后，二次项为

\[U(\alpha)=\frac12\left[mgl+
\frac{kq^2Ral}{D_0^2}\right]\alpha^2. \tag{24}\]

因此

\[\boxed{
f_4=2\pi\sqrt{
\frac{ml^2}{
mgl+\dfrac{kq^2R\left(lR+5l^2/4\right)}
{\left(l^2/16+lR/2\right)^2}}
}
}. \tag{25}\]

接地球半径趋于无穷大时应退化为接地平板，这是式（25）的重要核验。由

\[a=R+\frac54l\sim R,\]

以及

\[D_0=\frac{lR}{2}+\frac{l^2}{16}\sim\frac{lR}{2},\]

可得球面静电等效劲度项

\[\begin{aligned}
\frac{kq^2Ral}{D_0^2}
&\sim\frac{kq^2R\cdot R\cdot l}{(lR/2)^2}\\
&=\frac{4kq^2}{l}.
\end{aligned}\]

这恰好回到平板情形式（15）的静电劲度。因此

\[\lim_{R\to\infty}f_4=f_3.\]

另一方面令 \(q\to0\)，两个静电劲度项都消失，于是

\[f_3=f_4=2\pi\sqrt{\frac{l}{g}},\]

恢复普通单摆的小振动周期。两个极限分别检查了几何近似和电荷耦合。

前两问中要严格区分``质点所受合外力''和``绳的张力''。初态速度为零，所以径向加速度为零；沿绳方向的重力、电场力和张力必须恰好平衡，但沿切向仍有非零合力，质点随后才开始摆动。到最低点时速度已经由势能下降获得，径向方程中必须出现
\(mv_b^2/l\)，因此最低点张力不能用静力平衡直接求出。

平板镜像法给出的电场力为

\[F_e(h)=\frac{kq^2}{4h^2},\]

其中 \(2h\)
才是真电荷与像电荷的距离。若改用势能，接地体系的实际静电能是``真电荷与像电荷相互作用能''的一半，即
\(U_e=-kq^2/(4h)\)；这个 \(1/2\)
因子来自导体表面感应电荷建立过程中外界所做的功，不能省略。小振动时水平位移是
\(O(\alpha)\)，但摆球到导体的距离变化从 \(1-\cos\alpha=O(\alpha^2)\)
才开始，所以应把势能展开到 \(\alpha^2\)，再与
\(\tfrac12ml^2\dot\alpha^2\) 配对读取角频率。

接地球部分的像电荷大小和位置都随摆球到球心的距离 \(s\)
改变。最稳妥的做法是先写几何恒等式
\(s^2=a^2+l^2-2al\cos\alpha\)，再把接地球静电能写成只含 \(s\)
的函数，最后对 \(s(\alpha)\)
作二阶展开。这样可以避免分别展开像电荷量、像电荷位置和两电荷距离时产生相互抵消的繁琐项。令
\(q\to0\) 后两种周期都应退化为普通单摆周期
\(2\pi\sqrt{l/g}\)，这是必要的极限检查。

求初始合外力方向时，可先把竖直向下的 \(mg+F_e\)
分解到绳方向和切线方向。张力只调整径向分量，不能抵消切向分量，所以初始合外力沿切线指向最低点，其大小为
\((mg+F_e)\sin\alpha_0\)。若题目要求相对水平或竖直的方向，应把切线与竖直的夹角明确写出。最低点能量方程中的静电势能差可由
\(U_e(h)\) 直接计算，也可对位置相关电场力积分；两种方法应得到相同的
\(v_b^2\)。

镜像法中的``像电荷''只用于求真实区域的场：

\begin{enumerate}
\def\labelenumi{\arabic{enumi}.}
\tightlist
\item
  每个位置先标出真电荷、像电荷及二者距离，再由库仑力求真实电荷受力；接地平板能量若使用像电荷相互作用能，必须乘
  \(1/2\)。
\item
  摆动阶段用机械能求最低点速度，再在径向方程中求张力。能量方程和径向方程的力分量不能混写。
\item
  小振动应把重力势能和静电势能都展开到
  \(\alpha^2\)，常数项删除、一次项应在平衡点消失，二次项与
  \(ml^2\dot\alpha^2/2\) 配对读取频率。
\item
  接地球的像电荷位置与大小都随真电荷位置变；展开前先写精确几何量，最后用
  \(R\to\infty\) 检查是否退化为平板结果。
\end{enumerate}

\textbf{像电荷与势能：} 对距球心 \(a>R\) 的真电荷
\(q\)，接地球的像电荷参数为

\[q'=-q\frac Ra,\qquad a'=\frac{R^2}{a}.\]

平板则可视为曲率半径趋于无穷的极限；距板 \(d\) 时

\[F_z=-\frac{1}{4\pi\varepsilon_0}
\frac{q^2}{(2d)^2},\qquad
U=-\frac{1}{4\pi\varepsilon_0}\frac{q^2}{4d}.\]

能量中的 \(1/2\)
已包含在最后一式，不能再把真电荷---像电荷相互作用能原样照搬。

\subsection{六、关键洞察}\label{ux516dux5173ux952eux6d1eux5bdf-2}

镜像电荷可用于求力，但求接地体系的势能时必须乘
\(1/2\)，因为像电荷不是独立的真实电荷。小振动周期只取决于最低点附近势能的二阶项；静电力在最低点的数值本身不足以直接给出等效重力。

\subsection{七、易错点}\label{ux4e03ux6613ux9519ux70b9-2}

\begin{itemize}
\tightlist
\item
  把摆球与像电荷的距离写成 \(h\) 而不是 \(2h\)。
\item
  势能误写为 \(-kq^2/(2h)\)，遗漏接地体系的 \(1/2\) 因子。
\item
  初始张力中加入 \(mv^2/l\)；初始速度为零，径向加速度为零。
\item
  最低点只写受力方程而不先用能量求 \(v_b\)。
\item
  小振动中把 \(h-h_b\) 当作一阶小量；它实际从 \(\alpha^2\) 开始。
\item
  接地球情形把像电荷大小和位置当成常量。
\end{itemize}

\subsection{八、课堂讲解建议}\label{ux516bux8bfeux5802ux8bb2ux89e3ux5efaux8bae-2}

先独立回顾``平板镜像电荷的力''和``接地体系势能为何有一半''，再进入摆的动力学。小振动部分可让学生先判断距离变化是
\(O(\alpha)\) 还是
\(O(\alpha^2)\)，再展开势能。接地球部分重点讲清三个距离：悬点到球心
\(a\)、摆球到球心 \(s\)、最低点距离 \(s_0\)。

建议先把第（1）、（2）问分解为切向和径向两条线索。初态在径向上平衡，但切向合力不为零；最低点切向合力为零，但径向需要向心加速度。把这两个时刻并排比较，能够纠正学生把``速度为零''误解成``合力为零''、把``最低点''误解成``静力平衡''的常见错误。

讲势能时先用平板情形解释 \(1/2\) 因子，并要求学生从
\(F=-\mathrm dU/\mathrm dh\) 反向核验。小振动展开时把
\(1-\cos\alpha\approx\alpha^2/2\)
单独写出，说明导体距离的变化没有一阶项。接地球部分则先完成一张几何量表：\(a\)
是常量，\(s\) 随摆角变化，像电荷量和位置由 \(s\) 决定。统一写成
\(U_e(s)\) 后再展开，可显著降低符号负担。

可用两个极限收束课堂：\(q\to0\) 时必须回到普通单摆；\(R\to\infty\)
且保持最低点间距不变时，接地球局部应趋近接地平板。前一个适合全班完成，后一个可作为提高层的几何直觉检查。

\subsection{九、变式训练}\label{ux4e5dux53d8ux5f0fux8badux7ec3-2}

\begin{enumerate}
\def\labelenumi{\arabic{enumi}.}
\tightlist
\item
  把带电摆球换成电量为 \(-q\)，说明周期是否改变。
\item
  若接地平板改为固定电势但绝缘的均匀带电平板，比较其小振动周期与本题结果的结构差异。
\end{enumerate}

第一题训练识别镜像吸引力只依赖
\(q^2\)：真电荷变号时像电荷也随之变号，力仍为吸引，因此周期不变。第二题训练区分``接地导体的感应电荷''与``给定面电荷产生的外场''；后者若场强恒定，静电势能对高度为一次函数，其小振动效果等价于修改有效重力。

还可把初始角改为一般
\(\alpha_0\)，要求写出最低点张力的通式并判断绳不断裂所需的最大电荷量。该变式把能量、径向动力学和材料约束连接起来，解完后需检查
\(q\to0\) 时退化为普通单摆结论。

\textbf{变式解答：球面趋向平板。} 令 \(R\to\infty\)，有
\(a\sim R\)、\(D_0\sim Rl/2\)。接地球的电势能曲率系数趋于

\[\frac{kq^2Ral}{D_0^2}\longrightarrow\frac{4kq^2}{l},\]

故式（25）的周期趋于平板周期式（16）。这个极限同时检查镜像势能的 \(1/2\)
因子。若漏掉它，两个几何模型的结果会相差一倍静电刚度。

\newpage

\section{第 4 题}\label{ux7b2c-4-ux9898}

用单位长度电阻为 \(r\) 的均匀电阻丝围成图示矩形网络，各线段长度由
\(l\)、\(2l\)、\(4l\) 组成。灰色的 \(4l imes4l\)
中央区域内存在垂直纸面向里的均匀磁场

\[B=kt.\]

\begin{enumerate}
\def\labelenumi{\arabic{enumi}.}
\tightlist
\item
  求支路 \(AB\)、\(CD\) 中的电流。
\item
  按图中虚线接入理想电压表，求电压表读数；读数定义为图中``\(+\)''端电势减去``\(-\)''端电势。
\end{enumerate}

\begin{figure}[H]
\centering
\includegraphics[width=0.593\linewidth,height=\textheight,keepaspectratio,alt={图 4-1　时变磁场中的矩形电阻网络及电压表接法。}]{../figures/question-04.pdf}
\caption*{图 4-1　时变磁场中的矩形电阻网络及电压表接法。}
\end{figure}

\subsection{一、题型定位}\label{ux4e00ux9898ux578bux5b9aux4f4d-3}

时变磁场使闭合回路出现法拉第电动势。先确认九个网孔及各支路长度，再列网孔电流方程；开路电压还要指定引线的测量路径。

\subsection{二、知识点}\label{ux4e8cux77e5ux8bc6ux70b9-3}

\begin{itemize}
\tightlist
\item
  法拉第电磁感应定律和楞次定律。
\item
  均匀电阻丝的电阻 \(R=rL\)。
\item
  多网孔电流法，共用支路电流等于相邻网孔电流之差。
\item
  理想电压表不改变原网络电流。
\item
  时变磁场中两点间``电压''与选取的测量导线路径有关。
\end{itemize}

\subsection{三、物理图像}\label{ux4e09ux7269ux7406ux56feux50cf-3}

灰色区域的磁场垂直纸面向里且随时间增大，因此感应电场沿逆时针方向。不同网孔所围磁通变化率不同，电流不能凭几何对称直接判断。电压表两根引线与表头构成了指定的测量路径；沿这条路径计算的线积分才是表头读数。

\subsection{四、解题入口}\label{ux56dbux89e3ux9898ux5165ux53e3-3}

按图 4-2 规定四类由对称性合并的网孔电流 \(I_1,I_2,I_3,I_4\)
的正方向。每段电阻等于其长度乘 \(r\)。对每个独立回路写

\[\sum R(I_i-I_j)=\mathcal E_i,\]

其中 \(\mathcal E_i=-\mathrm d\Phi_i/\mathrm dt\)
的正负号由网孔正方向确定。求出网孔电流后，再由支路对应关系读取
\(I_{AB}\)、\(I_{CD}\)。

\subsection{五、解析}\label{ux4e94ux89e3ux6790-3}

图 4-2 明确规定四类由对称性合并的网孔电流 \(I_1,I_2,I_3,I_4\)
的绕向、磁场指向以及电压表正负端和虚线测量路径。网孔方程中的公共支路电流必须按这些绕向作差；求得网孔电流后，第（1）步读取
\(AB,CD\) 支路电流，第（2）步沿虚线路径计算 \(V_A-V_D\)。

\begin{figure}[H]
\centering
\includegraphics[width=0.591\linewidth,height=\textheight,keepaspectratio,alt={图 4-2　四类由对称性合并的网孔电流的规定方向与电压表测量路径。}]{../figures/solution-04.pdf}
\caption*{图 4-2　四类由对称性合并的网孔电流的规定方向与电压表测量路径。}
\end{figure}

灰色区域中

\[\frac{\mathrm dB}{\mathrm dt}=k.\]

若某回路包围的灰色面积为 \(A\)，其感应电动势大小为 \(Ak\)。

按照图示方向列四个网孔方程。把每一项中的公共长度单位 \(l\) 保留，得到

\[\begin{cases}
12lrI_1-4lrI_2-8lrI_4=8l^2k,\\
6lrI_2-2lrI_1-2lrI_3=0,\\
4lrI_3-lrI_2-lrI_4=0,\\
10lrI_4-2lrI_3-4lrI_1=4l^2k.
\end{cases} \tag{1}\]

第一式和第四式右端分别是相应回路所围磁通的变化率；第二、第三个网孔不包围净磁通，右端为零。

为看清代数过程，先把四式同除以 \(lr\)，并记

\[J=\frac{lk}{r}.\]

方程组化为

\[\begin{cases}
12I_1-4I_2-8I_4=8J,\\
-2I_1+6I_2-2I_3=0,\\
-I_2+4I_3-I_4=0,\\
-4I_1-2I_3+10I_4=4J.
\end{cases} \tag{2}\]

还可以把第一、第四式先各除以公因子，写成更贴近图形的形式：

\[3I_1-I_2-2I_4=2J,\]

\[5I_4-I_3-2I_1=2J.\]

第一式中的系数 \(3,1,2\)
分别对应本网孔总电阻长度以及与相邻网孔共用的边长倍数；第四式同理。右端同为
\(2J\)，表示除以电阻尺度后，两个围磁通网孔的有效感应电动势恰好相同。先写这种约分形式，再恢复到式（2），更容易把纯代数系数与图中的长度一一对应。

由第二、第三式分别得到

\[I_1=3I_2-I_3,\qquad I_4=4I_3-I_2. \tag{3}\]

把式（3）代入第一式：

\[12(3I_2-I_3)-4I_2-8(4I_3-I_2)=8J,\]

即

\[40I_2-44I_3=8J. \tag{4}\]

代入第四式：

\[-4(3I_2-I_3)-2I_3+10(4I_3-I_2)=4J,\]

即

\[-22I_2+42I_3=4J. \tag{5}\]

联立式（4）、（5）：

\[\begin{cases}
10I_2-11I_3=2J,\\
-11I_2+21I_3=2J.
\end{cases}\]

为避免把最后的分数结果当成黑箱，可继续做一次标准消元。第一式乘
\(11\)，第二式乘 \(10\)：

\[\begin{cases}
110I_2-121I_3=22J,\\
-110I_2+210I_3=20J.
\end{cases}\]

两式相加：

\[89I_3=42J,\]

所以

\[I_3=\frac{42}{89}J.\]

代回第一式：

\[10I_2=2J+11I_3
=2J+\frac{462}{89}J
=\frac{640}{89}J,\]

于是

\[I_2=\frac{64}{89}J.\]

因此消元结果可汇总为

\[I_2=\frac{64}{89}J,\qquad I_3=\frac{42}{89}J. \tag{6}\]

再由式（3）：

\[I_1=\frac{150}{89}J,\qquad I_4=\frac{104}{89}J. \tag{7}\]

这两个结果也可逐项还原：

\[\begin{aligned}
I_1
&=3I_2-I_3\\
&=3\cdot\frac{64}{89}J-\frac{42}{89}J\\
&=\frac{150}{89}J,
\end{aligned}\]

\[\begin{aligned}
I_4
&=4I_3-I_2\\
&=4\cdot\frac{42}{89}J-\frac{64}{89}J\\
&=\frac{104}{89}J.
\end{aligned}\]

把四个电流代回第二、第三个无感应电动势网孔，可得

\[-2I_1+6I_2-2I_3=0,\]

\[-I_2+4I_3-I_4=0,\]

说明公共边符号和代数求解彼此一致。

恢复 \(J=lk/r\)：

\[\boxed{
I_1=\frac{150lk}{89r},\quad
I_2=\frac{64lk}{89r},\quad
I_3=\frac{42lk}{89r},\quad
I_4=\frac{104lk}{89r}
}. \tag{8}\]

\textbf{（1）支路电流}

由图 4-2 可见，\(AB\) 与 \(CD\) 分别只属于网孔
\(I_2\)、\(I_4\)，因此它们的支路电流无需再作两个网孔电流之差；方向直接按图中箭头读取。

按图 4-2 的支路方向，\(AB\) 支路电流就是 \(I_2\)，\(CD\) 支路电流就是
\(I_4\)，故

\[\boxed{I_{AB}=\frac{64lk}{89r}},\qquad
\boxed{I_{CD}=\frac{104lk}{89r}}. \tag{9}\]

箭头方向以图 4-2 为准；若自行规定相反正方向，代数值应带负号。

\textbf{（2）电压表读数}

电压表的虚线路径穿过含感应电场的区域，不能只把沿途电阻压降机械相加。须按图
4-2 给定的正端 \(A\)、负端
\(D\)，沿同一条路径同时计入电阻电势降和感应电场线积分。

理想电压表内阻无穷大，不改变上述电流。表头正端接在 \(A\) 侧，负端接在
\(D\) 侧，读数定义为

\[U=V_A-V_D.\]

沿图中虚线所指定的测量路径，从 \(A\) 经 \(B,C\) 到 \(D\)
写电势降与感应电场线积分。按图示方向得到

\[2lrI_2+2lrI_3+4lrI_4-U=0. \tag{10}\]

因此

\[\begin{aligned}
U
&=2lr\frac{64lk}{89r}
+2lr\frac{42lk}{89r}
+4lr\frac{104lk}{89r}\\
&=\frac{128+84+416}{89}kl^2\\
&=\frac{628}{89}kl^2.
\end{aligned}\]

故电压表正向偏转，读数为

\[\boxed{U=V_A-V_D=\frac{628}{89}kl^2}. \tag{11}\]

代入过程还可按``每一段电势降''分开列账：

\[U_{AB}=2lrI_2=\frac{128}{89}kl^2,\]

\[U_{BC}=2lrI_3=\frac{84}{89}kl^2,\]

\[U_{CD}=4lrI_4=\frac{416}{89}kl^2.\]

于是

\[U=U_{AB}+U_{BC}+U_{CD}
=\frac{128+84+416}{89}kl^2.\]

三个分量都按图示测量方向为正，所以表头正端电势较高。若沿相反方向遍历同一条引线路径，三项与
\(U\) 会同时反号，物理读数的正负关系不会改变。

网孔法的关键不在解四元一次方程本身，而在每一项的来源。先统一规定四类由对称性合并的网孔电流的正方向，再对每个网孔沿同一绕行方向积分。某段只属于一个网孔时压降为
\(RI_i\)；某段被两个网孔共用时，实际支路电流是两个网孔电流的代数差，因此压降写成
\(R(I_i-I_j)\)。右端感应电动势的正负由楞次定律确定：灰区磁场向里增强，感应电流倾向产生向外磁场，所以逆时针绕灰区取正时电动势为正。

得到 \(I_1,I_2,I_3,I_4\)
后，还不能直接把某个网孔电流当成所求支路电流。应在图中标出支路两侧相邻网孔，并按所选支路正方向作代数相减。这样才能得到
\(I_{AB}\)、\(I_{CD}\)
的方向和大小。代回四个网孔方程逐式检验，是发现漏写公共电阻长度或面积系数的最快方法。

电压表问题还要多做一步路径说明。在时变磁场中
\(\nabla\times\boldsymbol E\ne0\)，不同引线路径可能围成不同磁通，所以``\(A\)
与 \(D\) 的电势差''不是脱离导线走向的唯一量。应沿题图虚线测量路径写

\[V_A-V_D=\int_D^A\boldsymbol E\cdot\mathrm d\boldsymbol l,\]

把经过的电阻压降和该闭合路径包围的感应电动势同时计入。最后用量纲检查：\(k\)
的单位为 \(\mathrm{T/s}\)，\(kl^2\) 的单位确为伏特；支路电流则应具有
\(lk/r\) 的量纲。

解线性方程组时不必追求口算。更重要的是保留每个系数的几何来源，并在求解后做三类检查：所有电流都与
\(k\) 成正比；把 \(k\)
反号时结果整体反号；任取一个不围磁通的回路，电阻压降代数和应为零。电压表理想开路，不改变原网络电流，这一点也应在列方程前说明。若把电压表支路误当成有电流的第五网孔，方程数量和物理模型都会错误。

网孔法的每个符号都应与图 4-2 的回路方向对应：

\begin{enumerate}
\def\labelenumi{\arabic{enumi}.}
\tightlist
\item
  对每个网孔先算磁通 \(\Phi=BA\)，再求
  \(\mathcal E=-\mathrm d\Phi/\mathrm dt\)；面积、绕向和公共支路方向共同决定右端符号。
\item
  公共电阻上的电流是相邻网孔电流的代数差。逐支路写电压降后再合并为四元线性方程，不能凭对称性猜电流。
\item
  解线性方程组时至少保留一次回代：先求公共分母，再把四个电流代回原方程检查每一行残差为零。
\item
  电压表读数沿指定路径逐段相加 \(IR\)
  电势降与感应电势；路径换向会整体反号，但读数大小应一致。
\end{enumerate}

\textbf{网孔电动势：} 四个网孔方程应统一写成线性系统

\[\boldsymbol A\boldsymbol I=\boldsymbol b,\qquad
\boldsymbol I=(I_1,I_2,I_3,I_4)^{\mathsf T}.\]

求解后必须同时做两类回代：

\[\begin{aligned}
\boldsymbol r&=\boldsymbol A\boldsymbol I-\boldsymbol b
=\boldsymbol0,\\
I_{AB}&=I_{\rm left}-I_{\rm right},\\
I_{CD}&=I_{\rm upper}-I_{\rm lower}.
\end{aligned}\]

第一行检查代数，后两行把计算变量还原为真正的支路电流，二者缺一不可。

\textbf{原图拓扑与系数对应。} 原题为三行三列共九个网孔，行高依次为
\(l,2l,l\)，列宽为 \(l,4l,l\)。图 4-2 中央灰区是
\(4l\times4l\)；不能漏掉底部的一条横线。反射对称后：中心网孔为
\(I_1\)，左右中间为 \(I_2\)，四角为 \(I_3\)，上下中间为
\(I_4\)，均规定逆时针。

中心周长为 \(12l\)，与左右各共用 \(2l\)，与上下各共用 \(4l\)，故第一式是
\(12lrI_1-4lrI_2-8lrI_4=8kl^2\)。上中央周长为 \(10l\)，与左右角各共用
\(l\)、与中心共用 \(4l\)，故第四式是
\(10lrI_4-2lrI_3-4lrI_1=4kl^2\)。非对称系数来自把等值网孔归为一类，不代表共用电阻的互易性失效。支路方向为
\(A\to B\)、\(C\to D\)。

\subsection{六、关键洞察}\label{ux516dux5173ux952eux6d1eux5bdf-3}

时变磁场中的电场不是保守场，因此``两个节点的电势差''不能脱离测量路径单独讨论。网孔方程中的感应电动势和电压表引线所围的磁通必须同时尊重图示方向。上述四元方程组正是本题最稳妥的主线。

\subsection{七、易错点}\label{ux4e03ux6613ux9519ux70b9-3}

\begin{itemize}
\tightlist
\item
  把灰色区域内的感应电动势平均分配到各边，导致与回路方程重复或漏算。
\item
  共用支路电流写成 \(I_i+I_j\)，没有根据箭头判断应为差还是和。
\item
  漏掉单位长度电阻 \(r\)，把长度直接当成电阻。
\item
  电压表正负端看反，只给读数绝对值。
\item
  把理想电压表接入后误认为网络中新增一条有电流的支路。
\item
  在非保守电场中随意更换从 \(A\) 到 \(D\) 的路径。
\end{itemize}

\subsection{八、课堂讲解建议}\label{ux516bux8bfeux5802ux8bb2ux89e3ux5efaux8bae-3}

先要求学生只做两件事：标出磁场增强时的感应电场方向；给每个网孔统一规定电流正方向。随后逐项解释第一条方程中的自电阻项、互电阻项和右端磁通项。电压表一问应明确写出
\(U=V_+-V_-\)，再沿指定导线走向累计电压。

板书第一条网孔方程时，应让学生逐段用笔沿回路走一遍：经过多长的电阻丝、该段是否与别的网孔共用、灰区面积是多少。不要一次写出四条方程后再解释，否则学生只会模仿系数。第一条完成后，让学生独立写第二条并互相检查公共边的符号。

求出网孔电流后回到题图，用不同颜色标出 \(AB\)、\(CD\)
两侧的网孔电流。强调``网孔电流是计算变量，支路电流才是物理可测量量''。可安排把答案代回原方程的快速检查，观察每条回路的电阻压降是否与感应电动势同量纲、同数量级。

电压表部分用两条不同引线做思想实验：若两种路径围住一块变化磁通，读数之差就等于这块磁通的变化率。这样学生能真正理解非保守感应电场中电压读数与引线布局有关，而不是把它记成一道特殊技巧题。

\subsection{九、变式训练}\label{ux4e5dux53d8ux5f0fux8badux7ec3-3}

\begin{enumerate}
\def\labelenumi{\arabic{enumi}.}
\tightlist
\item
  若磁场改为垂直纸面向外且
  \(B=kt\)，说明所有支路电流和电压表读数如何变化。
\item
  若电压表两根引线改从网络上方绕行，计算读数与本题有何不同，并用法拉第定律解释差值。
\end{enumerate}

第一题训练整体符号反演：\(\mathrm d\Phi/\mathrm dt\)
变号后，线性方程组右端整体变号，所有电流和给定路径的电压读数也整体反号，而大小不变。第二题训练测量路径意识；先把新旧引线合成一个闭合回路，再由法拉第定律直接求两次读数之差，比重新求完整网络更高效。

可进一步把灰区磁场改成左右两块、变化率分别为
\(k_1\)、\(k_2\)。让学生只改网孔方程右端，不重新推导电阻系数，从而体会线性电路中几何电阻矩阵与外部磁通源可以分开处理。

\textbf{变式解答：改变表头引线。}
若新旧测量路径围成的闭合回路按逆时针方向包围新增灰区面积
\(\Delta A\)，则按法拉第定律两条路径的线积分之差为
\(k\Delta A\)。统一端口和积分方向后，表头读数也相差相应的有符号
\(k\Delta A\)。引线只在零磁通区域内变形时读数不变。

\newpage

\section{第 5 题}\label{ux7b2c-5-ux9898}

水平光滑桌面上有两条平行光滑导轨，间距为
\(l\)。导轨所在区域存在垂直桌面向里的匀强磁场 \(B\)，左端用理想电感
\(L\) 连接。质量为 \(m\)
的导体棒横跨两导轨，并与一根``反向弹簧''相连：当棒相对 \(x=0\)
位置的位移为 \(x\) 时，弹簧给棒的力为 \(F=kx\)，方向指向远离 \(x=0\)
的一侧。忽略电阻以及弹簧的一切电磁效应。初始时 \(x=0\)，导体棒以速度
\(v_0\) 向右运动。

\begin{enumerate}
\def\labelenumi{\arabic{enumi}.}
\tightlist
\item
  初始电流为零。若反向弹簧始终保持理想特性，求使导体棒做匀速运动的劲度系数
  \(k\)。
\item
  若 \(k\) 取第（1）问结果的一半，求导体棒的振动周期和平衡位置。
\item
  在第（2）问条件下，求导体棒位移的最大值 \(x_{3\max}\)、最小值
  \(x_{3\min}\)，以及电流绝对值的最大值 \(I_{3m}\)。
\item
  仍取第（1）问劲度系数的一半，但初始时电路中已有大小为 \(I_0\)
  的逆时针电流。把 \(I_0\) 视为小量，求位移最大值 \(x_{4\max}\) 和最小值
  \(x_{4\min}\)，并给出一阶小量结果。
\end{enumerate}

\begin{figure}[H]
\centering
\includegraphics[width=0.9\linewidth,height=\textheight,keepaspectratio,alt={图 5-1　磁场中由理想电感闭合的导轨---导体棒系统。}]{../figures/question-05.pdf}
\caption*{图 5-1　磁场中由理想电感闭合的导轨---导体棒系统。}
\end{figure}

\subsection{一、题型定位}\label{ux4e00ux9898ux578bux5b9aux4f4d-4}

理想电感将导体棒位移与电流联系起来，安培力形成回复作用。它与反向弹簧竞争，决定棒是匀速、简谐振动还是失稳运动。

\subsection{二、知识点}\label{ux4e8cux77e5ux8bc6ux70b9-4}

\begin{itemize}
\tightlist
\item
  导体棒的动生电动势 \(\mathcal E=Bl\dot x\)。
\item
  理想电感关系 \(L\dot I=\mathcal E\)。
\item
  安培力 \(F_B=BlI\) 及楞次定律确定方向。
\item
  对时间积分得到电流---位移约束。
\item
  平衡位置平移和简谐振动振幅。
\item
  小参数展开与一阶修正。
\end{itemize}

\subsection{三、物理图像}\label{ux4e09ux7269ux7406ux56feux50cf-4}

导体棒向右运动时，回路所围面积和向里的磁通增加，感应电流为逆时针，棒中电流向上，安培力向左。由于没有电阻，电流不会瞬时耗散；从初态积分后，电流与棒的位移成正比。因此磁场和电感合起来产生一个等效劲度系数
\(B^2l^2/L\) 的回复力。

初始已有电流时，电流---位移关系多出常量
\(I_0\)，安培力也多出恒定项。它不改变振动频率，只把平衡位置整体向左平移。

\subsection{四、解题入口}\label{ux56dbux89e3ux9898ux5165ux53e3-4}

规定 \(x\) 向右为正，逆时针电流为正。写出

\[L\frac{\mathrm dI}{\mathrm dt}=Bl\dot x,\qquad
m\ddot x=kx-BlI. \tag{1}\]

第一式对时间积分，不要先对第二式求导。这样可以直接把 \(I\) 用 \(x\)
表示，系统立即化为一个二阶常系数方程。

\subsection{五、解析}\label{ux4e94ux89e3ux6790-4}

图 5-2 用等效势 \(U_{\rm eff}(x)\)
表示弹簧力与磁场---电感恢复力的合成：初始电流 \(I_0=0\) 时平衡点在
\(x=0\)，\(I_0>0\) 时抛物线整体平移至
\(x_e=-2LI_0/(Bl)\)；橙色箭头标出安培力方向。第（1）---（3）步处理零初始电流，第（4）步使用平移后的平衡点。

\begin{figure}[H]
\centering
\includegraphics[width=0.739\linewidth,height=\textheight,keepaspectratio,alt={图 5-2　动生电流、安培力方向及等效振子平衡位置。}]{../figures/solution-05.pdf}
\caption*{图 5-2　动生电流、安培力方向及等效振子平衡位置。}
\end{figure}

\textbf{（1）匀速运动条件}

图 5-2
所示磁场---电感效应产生与位移成正比的恢复力。先把法拉第---电感方程积分得到
\(I(x)\)，再与反向弹簧力比较；只有二者系数恰好抵消时，初速度才保持不变。

初始条件为

\[x(0)=0,\qquad \dot x(0)=v_0,\qquad I(0)=0.\]

对式（1）的电感方程积分：

\[L\int_0^I\mathrm dI=Bl\int_0^x\mathrm dx,\]

得到

\[LI=Blx,\qquad I=\frac{Bl}{L}x. \tag{2}\]

代入机械方程：

\[m\ddot x=kx-Bl\left(\frac{Bl}{L}x\right)
=\left(k-\frac{B^2l^2}{L}\right)x. \tag{3}\]

要使棒保持初始速度做匀速直线运动，需要任意位移下合力均为零，即式（3）中
\(x\) 的系数为零：

\[\boxed{k_0=\frac{B^2l^2}{L}}. \tag{4}\]

此时 \(\ddot x=0\)，所以
\(x=v_0t\)。电流仍随位移线性增加，但其安培力始终恰好抵消反向弹簧力。

\textbf{（2）劲度系数减半}

取

\[k=\frac12k_0=\frac{B^2l^2}{2L}. \tag{5}\]

式（3）变为

\[m\ddot x=-\frac{B^2l^2}{2L}x,\]

即

\[\ddot x+\omega^2x=0,\qquad
\omega^2=\frac{B^2l^2}{2mL}. \tag{6}\]

平衡位置是

\[\boxed{x_{\mathrm e}=0}. \tag{7}\]

周期为

\[\boxed{
T=\frac{2\pi}{\omega}
=2\pi\sqrt{\frac{2mL}{B^2l^2}}
}. \tag{8}\]

\textbf{（3）位移极值与最大电流}

初始条件 \(x(0)=0\)、\(\dot x(0)=v_0\) 给出

\[x(t)=\frac{v_0}{\omega}\sin\omega t. \tag{9}\]

振幅

\[A=\frac{v_0}{\omega}
=v_0\sqrt{\frac{2mL}{B^2l^2}}. \tag{10}\]

因此

\[\boxed{x_{3\max}=v_0\sqrt{\frac{2mL}{B^2l^2}}},\]

\[\boxed{x_{3\min}=-v_0\sqrt{\frac{2mL}{B^2l^2}}}. \tag{11}\]

由式（2），电流绝对值在 \(|x|=A\) 时最大：

\[\begin{aligned}
I_{3m}
&=\frac{Bl}{L}A\\
&=\frac{Bl}{L}v_0\sqrt{\frac{2mL}{B^2l^2}}\\
&=v_0\sqrt{\frac{2m}{L}}.
\end{aligned}\]

故

\[\boxed{I_{3m}=v_0\sqrt{\frac{2m}{L}}}. \tag{12}\]

\textbf{（4）初始存在逆时针电流 \(I_0\)}

有初始电流时，图 5-2 红色等效势的最低点移到 \(x_e<0\)。这来自积分常量
\(LI_0\) 产生的恒定安培力项；必须先求新平衡位置，再围绕 \(x_e\)
写振动，不能仍把 \(x=0\) 当作平衡点。

此时积分常量不能丢掉。由电感方程

\[L[I(t)-I_0]=Bl[x(t)-0],\]

所以

\[I=\frac{Bl}{L}x+I_0. \tag{13}\]

仍取 \(k=B^2l^2/(2L)\)。代入机械方程：

\[\begin{aligned}
m\ddot x
&=\frac{B^2l^2}{2L}x
-Bl\left(\frac{Bl}{L}x+I_0\right)\\
&=-\frac{B^2l^2}{2L}x-BlI_0.
\end{aligned}\]

整理为

\[\ddot x+\omega^2x=-\frac{Bl}{m}I_0, \tag{14}\]

其中 \(\omega\) 仍由式（6）给出。令右端对应的平衡位置为
\(x_{\mathrm e}\)：

\[\omega^2x_{\mathrm e}=-\frac{Bl}{m}I_0.\]

于是

\[\boxed{x_{\mathrm e}=-\frac{2LI_0}{Bl}}. \tag{15}\]

引入相对平衡位置的坐标

\[y=x-x_{\mathrm e}=x+\frac{2LI_0}{Bl}.\]

则

\[\ddot y+\omega^2y=0. \tag{16}\]

初始条件变为

\[y(0)=\frac{2LI_0}{Bl},\qquad \dot y(0)=v_0.\]

相对平衡位置的振幅为

\[\begin{aligned}
A_4
&=\sqrt{y^2(0)+\frac{\dot y^2(0)}{\omega^2}}\\
&=\sqrt{\frac{4L^2I_0^2}{B^2l^2}
+\frac{2mLv_0^2}{B^2l^2}}.
\end{aligned} \tag{17}\]

因此精确的两个极值为

\[\boxed{
x_{4\max}
=-\frac{2LI_0}{Bl}
+\sqrt{\frac{4L^2I_0^2+2mLv_0^2}{B^2l^2}}
}, \tag{18}\]

\[\boxed{
x_{4\min}
=-\frac{2LI_0}{Bl}
-\sqrt{\frac{4L^2I_0^2+2mLv_0^2}{B^2l^2}}
}. \tag{19}\]

若 \(I_0\) 是一阶小量，则根号中 \(I_0\)
以平方出现，只贡献二阶修正。因此保留到 \(O(I_0)\)：

\[\boxed{
x_{4\max}\approx v_0\sqrt{\frac{2mL}{B^2l^2}}
-\frac{2LI_0}{Bl}
}, \tag{20}\]

\[\boxed{
x_{4\min}\approx- v_0\sqrt{\frac{2mL}{B^2l^2}}
-\frac{2LI_0}{Bl}
}. \tag{21}\]

本题第一步必须从完整回路方程出发。按图中规定的电流正方向，导体棒位移改变回路磁通，法拉第定律给出
\(L\dot I-Bl\dot x=0\)；对时间积分后得到

\[LI-Blx=C.\]

积分常量由初态决定。初始电流为零时 \(C=0\)，所以 \(I=Blx/L\)；有初始电流
\(I_0\) 时则为 \(C=LI_0\)，即
\(I=Blx/L+I_0\)。这正是第（4）问不能直接沿用前三问电流---位移关系的原因。

安培力方向应由楞次定律而不是凭公式猜测：棒向右运动使向里的磁通增加，感应电流产生向外磁场，棒受力向左，因此磁力是回复力。把
\(I(x)\) 代入 \(F_B=IlB\)，再与反向弹簧力 \(+kx\) 合并，可得等效系数
\(B^2l^2/L-k\)。当它为零时合力恒为零，棒做匀速运动；取临界劲度的一半后，合力与位移反向且正比，才形成简谐振动。

有 \(I_0\) 时方程多出常力项。先配方写成关于 \(x-x_e\)
的标准振子，比直接求通解更清楚：固有频率由二次项决定，初始电流只平移平衡位置。极值由``相对平衡位置的振幅''给出，根号内的
\(I_0^2\) 是二阶量，而平衡位置平移是一次量，所以一阶近似只保留后者。令
\(I_0\to0\) 应退化到第（3）问，令 \(v_0\to0\)
时两个极值应以新平衡点为中心对称。

还可用能量作独立核验。系统守恒量为

\[E=\frac12m\dot x^2+\frac12LI^2-\frac12kx^2,\]

其中反向弹簧的势能带负号。把 \(LI-Blx=LI_0\) 代入后，\(E\) 对 \(x\)
成为开口向上的二次函数，当 \(k<k_0\) 时存在稳定最低点；\(k=k_0\)
时二次项消失；\(k>k_0\)
时能量沿某方向无下界。这个能量图像与微分方程的振动、匀速和失稳三种结论完全一致。

电感回路的积分常量决定第（4）问的平衡位置：

\begin{enumerate}
\def\labelenumi{\arabic{enumi}.}
\tightlist
\item
  法拉第---电感方程先积分为 \(L[I(t)-I(0)]=Bl[x(t)-x(0)]\)。只有
  \(I(0)=x(0)=0\) 时才能简写为 \(LI=Blx\)。
\item
  按图 5-2
  用右手定则确定安培力总与引起电流的运动相反，再与题设反向弹簧力一起写入牛顿方程。
\item
  合并 \(x\) 的系数后识别 \(\omega^2\)。若有常力项，先配方求真实平衡位置
  \(x_e\)，再令 \(y=x-x_e\)，不能绕原点直接写简谐振动。
\item
  极值由振幅给出，最大电流再由含初始常量的 \(I(x)\)
  求；最后用机械能加电感能守恒检查两个端点。
\end{enumerate}

\textbf{电流与棒的位移：} 电感方程必须先积分，再代入机械方程：

\[\begin{aligned}
L\frac{\mathrm dI}{\mathrm dt}&=Bl\frac{\mathrm dx}{\mathrm dt},\\
LI-Blx&=LI_0-Blx_0,\\
m\ddot x&=kx-BlI(x).
\end{aligned}\]

整理为 \(m\ddot x+k_{\rm eff}x=F_{\rm eff}\) 后，应先求

\[x_e=\frac{F_{\rm eff}}{k_{\rm eff}},\qquad
y=x-x_e,\]

再对 \(y\) 写简谐振动；否则振幅和最大电流都会少掉积分常量。

\subsection{六、关键洞察}\label{ux516dux5173ux952eux6d1eux5bdf-4}

理想电感使电流记录了导体棒从初态开始积累的位移，因此磁力是``位移型回复力''而不是``速度型阻力''。初始电流只改变积分常量，相当于给振子增加恒力并平移平衡位置；固有频率不变。

\subsection{七、易错点}\label{ux4e03ux6613ux9519ux70b9-4}

\begin{itemize}
\tightlist
\item
  把电感当成电阻，写成 \(I=\mathcal E/L\)。
\item
  忽略积分常量，导致第（4）问仍写 \(LI=Blx\)。
\item
  安培力方向写错，使磁力和反向弹簧力同向，错误地判断系统必然失稳。
\item
  第（2）问把平衡位置误认为初始速度决定的非零位置。
\item
  求最大电流时只取正位移端点而不说明取的是绝对值。
\item
  小量展开中保留根号里的 \(I_0^2\) 却遗漏平衡位置的一阶平移。
\end{itemize}

\subsection{八、课堂讲解建议}\label{ux516bux8bfeux5802ux8bb2ux89e3ux5efaux8bae-4}

先用楞次定律确定棒向右时电流和安培力方向，再写符号方程。讲完式（2）后可暂停，让学生解释为什么没有电阻却仍能出现回复力。第（4）问建议先画出``原平衡点''和``新平衡点''，再求相对坐标振幅。

建议先只研究电磁子系统。由 \(L\dot I-Bl\dot x=0\) 积分得到 \(LI-Blx=C\)
后，让学生解释电流为什么``记住的是位移而不是瞬时速度''。随后再接入机械方程，这样可以避免把安培力误当成与速度成正比的阻尼力。方向判断应至少用楞次定律口头核验一次。

第（1）问把三种情况放在同一张相图上：\(k<k_0\) 为稳定振动，\(k=k_0\)
为匀速，\(k>k_0\)
为指数失稳。第（2）、（3）问统一改写为标准简谐振子，再从初始位置和速度读振幅。第（4）问先把常力项移到等号右边，画出平衡点从
\(0\) 移到 \(x_e\)，再在新坐标 \(X=x-x_e\) 中处理。

一阶近似是很好的数学训练点。要求学生明确写出 \(I_0\) 的阶数：根号内的
\(I_0^2\) 可舍去，但平衡位置中的 \(I_0\) 必须保留。最后用 \(I_0\to0\) 和
\(v_0\to0\) 两个极限检查极值表达式。

\subsection{九、变式训练}\label{ux4e5dux53d8ux5f0fux8badux7ec3-4}

\begin{enumerate}
\def\labelenumi{\arabic{enumi}.}
\tightlist
\item
  若回路中再串联一个电阻
  \(R_0\)，写出耦合微分方程，并判断磁力是否仍能简单等效为弹簧。
\item
  若 \(k>k_0\)，求位移随时间的形式并解释系统为何指数失稳。
\end{enumerate}

第一题训练从守恒耦合系统过渡到耗散系统。加入电阻后，电流不再是位移的代数函数，而要与
\(x\)
联立求解；消元后出现速度相关项，机械能通过焦耳热减少。第二题训练辨认负的等效劲度系数，通解由三角函数变为双曲函数或指数函数，初始微小偏离会被放大。

还可令磁场随位置缓慢变化 \(B(x)=B_0(1+\varepsilon x/l)\)，只保留
\(\varepsilon\) 的一阶项。该变式要求学生重新积分磁通而不能直接把常量
\(B\) 代入原公式，并可观察非线性回复力怎样产生振幅相关的周期修正。

\textbf{变式解答：一般劲度。} 令 \(K=B^2l^2/L-k\)。若
\(K>0\)，\(\omega=\sqrt{K/m}\)、\(x_e=-BlI_0/K\)。初态
\(x=0\)、\(\dot x=v_0\) 给振幅 \(A=\sqrt{x_e^2+v_0^2/\omega^2}\)，极值为
\(x_e\pm A\)。若 \(K=0\) 且 \(I_0=0\)，运动匀速；若 \(K=0\) 但
\(I_0\ne0\)，为常加速度运动。\(K<0\)
时齐次解是指数函数，不能再称为振动。

\newpage

\section{第 6 题}\label{ux7b2c-6-ux9898}

把原恒星和现在的太阳都近似为均匀球体。原恒星早期半径远大于现在的太阳半径，可视为无穷大。已知太阳表面温度
\(T=6000\,\mathrm K\)、质量 \(M=2.0\times10^{30}\,\mathrm{kg}\)、半径
\(R=7.0\times10^8\,\mathrm m\)，万有引力常量
\(G=6.7\times10^{-11}\,\mathrm{N\,m^2\,kg^{-2}}\)，斯特藩---玻尔兹曼常量
\(\sigma=5.7\times10^{-8}\,\mathrm{W\,m^{-2}\,K^{-4}}\)。

\begin{enumerate}
\def\labelenumi{\arabic{enumi}.}
\tightlist
\item
  估算太阳由原恒星收缩到现在半径时释放的引力势能。
\item
  若按太阳当前的黑体辐射功率辐射，这些能量相当于太阳多长时间的辐射能量？
\item
  若第（1）问的能量瞬间全部转化为太阳内能，取太阳平均比热容
  \(c=10^3\,\mathrm{J\,kg^{-1}\,K^{-1}}\)，估算可使太阳达到的温度。氢聚变点火温度约为
  \(T_C=2.0\times10^7\,\mathrm K\)。假设一般恒星的平均密度和比热容与太阳相同，估算能够达到点火温度的恒星质量至少为太阳质量的百分之多少。
\end{enumerate}

\subsection{一、题型定位}\label{ux4e00ux9898ux578bux5b9aux4f4d-5}

按题设计算自引力球收缩释放的能量，再以黑体辐射功率估算寿命和点火温度。能量转化比例是模型条件，需与维里平衡情形区分。

\subsection{二、知识点}\label{ux4e8cux77e5ux8bc6ux70b9-5}

\begin{itemize}
\tightlist
\item
  均匀球体的引力自能 \(U=-3GM^2/(5R)\)。
\item
  斯特藩---玻尔兹曼定律 \(P=4\pi R^2\sigma T^4\)。
\item
  热量估算 \(Q=Mc\Delta T\)。
\item
  同密度恒星的质量---半径关系 \(R\propto M^{1/3}\)。
\item
  数量级计算、单位换算和比例标度。
\end{itemize}

\subsection{三、物理图像}\label{ux4e09ux7269ux7406ux56feux50cf-5}

把均匀球体看成从内到外逐层装配。将一层质量 \(\mathrm dm\)
从无穷远搬到半径 \(r\) 处，释放的引力势能大小为

\[\frac{Gm(r)\,\mathrm dm}{r}.\]

对所有球壳积分，就得到均匀球的总引力自能。

恒星质量越大，引力势阱越深。若平均密度相同，质量增加时半径也增加，但只按
\(M^{1/3}\) 增加，因此单位质量可获得的引力能按 \(GM/R\propto M^{2/3}\)
增长。

\subsection{四、解题入口}\label{ux56dbux89e3ux9898ux5165ux53e3-5}

不要直接把全部质量当成位于球心的点质量。先推导均匀球的 \(3/5\)
系数，再代数值。第三问先求太阳自身若完全保留收缩能量可达到的温度
\(T_s\)，随后用

\[\frac{T'}{T_s}=\left(\frac{M'}{M}\right)^{2/3}\]

求点火质量。

\subsection{五、解析}\label{ux4e94ux89e3ux6790-5}

图 6-1 左侧把均匀球看成从无穷远逐层搬入薄壳
\(dm=4\pi r^2\rho\,dr\)，内部已有质量为 \(m(r)\)；右侧把所得结合能
\(E_g=3GM^2/(5R)\)
分成``按黑体功率辐射''与``转化为内能升温''两条能量路径。三问分别沿这条能量链计算自能、辐射时间和点火质量。

\begin{figure}[H]
\centering
\includegraphics[width=0.9\linewidth,height=\textheight,keepaspectratio,alt={图 6-1　均匀球薄壳装配，以及收缩能量向辐射寿命和恒星升温的转换关系。}]{../figures/solution-06.pdf}
\caption*{图
6-1　均匀球薄壳装配，以及收缩能量向辐射寿命和恒星升温的转换关系。}
\end{figure}

\textbf{（1）太阳收缩释放的引力势能}

沿图 6-1 左侧装配半径为 \(r\) 的薄壳时，只有壳与内部质量 \(m(r)\)
的相互作用是本次新增势能，故 \(dU=-Gm(r)dm/r\)；从 \(r=0\) 积到 \(R\)
可避免重复计算壳层之间的相互作用。

设太阳平均密度为 \(\rho\)。半径 \(r\) 内的质量

\[m(r)=\frac43\pi r^3\rho,\]

厚度为 \(\mathrm dr\) 的球壳质量

\[\mathrm dm=4\pi r^2\rho\,\mathrm dr.\]

把该球壳从无穷远搬到 \(r\) 处，引力势能变化为

\[\mathrm dU=-\frac{Gm(r)\mathrm dm}{r}. \tag{1}\]

代入 \(m(r)\) 和 \(\mathrm dm\)：

\[\mathrm dU
=-\frac{G}{r}\left(\frac43\pi r^3\rho\right)
\left(4\pi r^2\rho\,\mathrm dr\right)
=-\frac{16}{3}\pi^2G\rho^2r^4\,\mathrm dr. \tag{2}\]

从 \(0\) 积分到 \(R\)：

\[U=-\frac{16}{3}\pi^2G\rho^2\int_0^Rr^4\,\mathrm dr
=-\frac{16}{15}\pi^2G\rho^2R^5. \tag{3}\]

又

\[M=\frac43\pi R^3\rho
\quad\Rightarrow\quad
\rho=\frac{3M}{4\pi R^3}.\]

代入式（3）：

\[U=-\frac35\frac{GM^2}{R}. \tag{4}\]

原恒星早期半径视为无穷大，引力自能近似为零，所以收缩释放的能量大小为

\[E_g=\frac{3GM^2}{5R}. \tag{5}\]

代入数值：

\[\begin{aligned}
E_g
&=\frac{3\times6.7\times10^{-11}
\times(2.0\times10^{30})^2}
{5\times7.0\times10^8}\\
&\approx2.30\times10^{41}\,\mathrm J.
\end{aligned}\]

因此

\[\boxed{E_g\approx2.3\times10^{41}\,\mathrm J}. \tag{6}\]

量纲上

\[\left[\frac{GM^2}{R}\right]
=\frac{\mathrm{N\,m^2\,kg^{-2}}\cdot\mathrm{kg^2}}{\mathrm m}
=\mathrm{N\,m}=\mathrm J,\]

与能量一致。若把原恒星半径记为有限的 \(R_i\)，更一般的释放能量应为

\[\Delta E_g=\frac{3GM^2}{5}
\left(\frac1R-\frac1{R_i}\right).\]

本题使用 \(R_i\gg R\)，所以 \(1/R_i\) 可忽略，才得到式（5）。

\textbf{（2）按当前辐射功率对应的时间}

图 6-1 右侧的辐射支路先由斯忒藩---玻尔兹曼定律求总功率
\(P=4\pi R^2\sigma T^4\)，再用能量除以功率。这里假定整个收缩能量以当前恒定光度辐射，只是量级估算。

太阳表面积为 \(4\pi R^2\)。按黑体辐射定律，辐射功率

\[P=4\pi R^2\sigma T^4. \tag{7}\]

逐项估算：

\[T^4=(6000)^4=1.296\times10^{15}\,\mathrm K^4,\]

\[4\pi R^2=4\pi(7.0\times10^8)^2
\approx6.16\times10^{18}\,\mathrm{m^2}.\]

故

\[P\approx6.16\times10^{18}
\times5.7\times10^{-8}
\times1.296\times10^{15}
\approx4.55\times10^{26}\,\mathrm W. \tag{8}\]

相应时间

\[t=\frac{E_g}{P}
\approx\frac{2.30\times10^{41}}{4.55\times10^{26}}
\approx5.05\times10^{14}\,\mathrm s. \tag{9}\]

用 \(1\,\mathrm{yr}\approx3.16\times10^7\,\mathrm s\) 换算：

\[t\approx1.60\times10^7\,\mathrm{yr}. \tag{10}\]

所以

\[\boxed{t\approx5.1\times10^{14}\,\mathrm s
\approx1.6\times10^7\,\text{年}}.\]

秒到年的换算完整写为

\[\begin{aligned}
t
&=\frac{5.05\times10^{14}\,\mathrm s}
{3.16\times10^7\,\mathrm{s/yr}}\\
&=1.60\times10^7\,\mathrm{yr}.
\end{aligned}\]

与 \(5\times10^9\,\mathrm{yr}\) 的太阳年龄比较：

\[\frac{5\times10^9}{1.6\times10^7}\approx3.1\times10^2.\]

也就是说，引力收缩按当前光度只能维持太阳年龄的大约三百分之一，数量级矛盾非常显著。

这远小于太阳约 \(5\times10^9\) 年的年龄。

\textbf{（3）引力升温与最小点火质量}

沿图 6-1 右侧的升温支路，令结合能等于热容量所需能量 \(McT_s\)，先得到
\(T_s\propto M/R\)；再结合相同平均密度下 \(R\propto M^{1/3}\)，得到
\(T_s\propto M^{2/3}\) 并反求临界质量。

若全部收缩能量转化为太阳内能，并忽略初温，则

\[McT_s=E_g=\frac{3GM^2}{5R}.\]

约去一个 \(M\)：

\[T_s=\frac{3GM}{5Rc}. \tag{11}\]

代入数值：

\[\begin{aligned}
T_s
&=\frac{3\times6.7\times10^{-11}\times2.0\times10^{30}}
{5\times7.0\times10^8\times10^3}\\
&\approx1.15\times10^8\,\mathrm K.
\end{aligned}\]

故

\[\boxed{T_s\approx1.15\times10^8\,\mathrm K}. \tag{12}\]

再设另一颗恒星的质量、半径为 \(M',R'\)。平均密度与太阳相同意味着

\[\frac{M'}{R'^3}=\frac{M}{R^3},\]

所以

\[\frac{R'}{R}=\left(\frac{M'}{M}\right)^{1/3}. \tag{13}\]

其引力升温为

\[T' =\frac{3GM'}{5R'c}.\]

与太阳的式（11）相除：

\[\frac{T'}{T_s}
=\frac{M'/M}{R'/R}
=\left(\frac{M'}{M}\right)^{2/3}. \tag{14}\]

令 \(T'=T_C\)，得到

\[\frac{M'}{M}=\left(\frac{T_C}{T_s}\right)^{3/2}. \tag{15}\]

代入 \(T_C=2.0\times10^7\,\mathrm K\)：

\[\frac{M'}{M}
=\left(\frac{2.0\times10^7}{1.15\times10^8}\right)^{3/2}
\approx0.073. \tag{16}\]

指数 \(3/2\) 的数值计算可拆开：

\[\frac{T_C}{T_s}\approx0.174,\]

\[\left(0.174\right)^{3/2}
=0.174\sqrt{0.174}
\approx0.174\times0.417
\approx0.073.\]

换成百分数：

\[0.073\times100\%=7.3\%.\]

由于 \(T'\propto M'^{2/3}\) 单调增加，满足点火条件的是
\(M'\geq M'_{\min}\)，所以该结果应解释为质量下限。

因此

\[\boxed{M'_{\min}\approx7.3\%\,M}. \tag{17}\]

引力自能应通过``逐层装配''得到。把半径 \(r\) 内已有质量记为
\(M(r)=M r^3/R^3\)，再把质量为 \(\mathrm dm=4\pi r^2\rho\,\mathrm dr\)
的薄壳从无穷远搬到 \(r\) 处，外界释放的微元能量为
\(G M(r)\mathrm dm/r\)。从 \(0\) 积分到 \(R\) 才得到系数
\(3/5\)；若只用数量级式 \(GM^2/R\)，会丢失题目要求的数值系数。

第二问的辐射功率来自球面面积乘黑体单位面积辐射功率：\(P=4\pi R^2\sigma T^4\)。先分别计算
\(R^2\) 和 \(T^4\)，再组合指数，可减少科学记数法错误。时间 \(t=E/P\)
的单位是秒，换算成年时还要除以 \(3.15\times10^7\)。所得约
\(1.6\times10^7\)
年远小于太阳年龄，这个矛盾正是历史上引出核能源机制的物理意义。

第三问严格按题设把收缩能全部转化为内能，所以写
\(McT_s=E_g\)；这里不要擅自再乘维里定理的
\(1/2\)。推广到一般恒星时，同平均密度意味着
\(R\propto M^{1/3}\)，于是单位质量获得的引力能
\(GM/R\propto M^{2/3}\)，温度也按同样幂次标度。求得质量下限后要检查趋势：质量越大，\(M/R\)
越大，内部越容易达到点火温度，因此临界值应是下限而不是上限。

数值代入可分层记录有效数字。引力能约为
\(2.3\times10^{41}\,\mathrm J\)；太阳光度由给定数据算得约
\(4.6\times10^{26}\,\mathrm W\)，两者相除约
\(5.0\times10^{14}\,\mathrm s\)。温度估算 \(E_g/(Mc)\)
中质量约掉一个数量级后得到
\(10^8\,\mathrm K\)，应高于给定点火温度。若某一步得到
\(10^5\,\mathrm K\) 或 \(10^{11}\)
年，应首先检查半径、温度四次方和秒---年的换算指数。

图 6-1 把``逐层装配---辐射---升温''三段能量链放在同一视图中：

\begin{enumerate}
\def\labelenumi{\arabic{enumi}.}
\tightlist
\item
  薄壳积分中，半径 \(r\) 处只能使用已经装好的内部质量 \(m(r)\)。将
  \(m(r)\) 与 \(\mathrm dm\) 全部代入后再积分，系数 \(3/5\)
  才有明确来源。
\item
  黑体功率按``表面积乘单位面积功率''计算。\(R^2\)、\(T^4\) 和 \(\sigma\)
  的十次幂先分别计算，再相乘；能量除以功率得到秒，最后显式除以每年的秒数。
\item
  升温方程按题设写 \(McT_s=E_g\)，约去一个 \(M\)
  后再代数值。这里不能擅自加入未在题设中的维里定理系数。
\item
  同密度给出 \(R'/R=(M'/M)^{1/3}\)，再由 \(T\propto M/R\) 得
  \(T\propto M^{2/3}\)。解出临界质量后要检查质量越大温度越高，所以结果是下限。
\end{enumerate}

\textbf{收缩能量积分：} 均匀球逐壳装配时

\[\begin{aligned}
m(r)&=\frac{4\pi}{3}\rho r^3,
&\mathrm dm&=4\pi\rho r^2\,\mathrm dr,\\
E_g&=\int_0^R\frac{Gm(r)}{r}\,\mathrm dm\\
&=\frac{16\pi^2G\rho^2}{3}\int_0^Rr^4\,\mathrm dr
=\frac{3GM^2}{5R}.
\end{aligned}\]

这里的 \(E_g\) 是释放的结合能大小；若改写为系统引力势能，则应记为
\(U_g=-E_g\)。

\subsection{六、关键洞察}\label{ux516dux5173ux952eux6d1eux5bdf-5}

均匀球引力自能中的 \(3/5\) 来自逐层装配，而不是简单的
\(GM^2/R\)。同密度条件把温度标度变为
\(T\propto M^{2/3}\)，所以点火质量对温度比取 \(3/2\)
次方。第二问得到的约 \(10^7\) 年就是太阳的开尔文---亥姆霍兹时间数量级。

\subsection{七、易错点}\label{ux4e03ux6613ux9519ux70b9-5}

\begin{itemize}
\tightlist
\item
  把引力自能写成 \(GM^2/R\)，遗漏 \(3/5\)。
\item
  在装配球体时用总质量 \(M\) 代替半径 \(r\) 内已经装好的质量 \(m(r)\)。
\item
  黑体辐射功率遗漏表面积 \(4\pi R^2\)。
\item
  把 \(6000^4\) 的数量级算成 \(10^{12}\)。
\item
  第三问写成 \(McT=3GM^2/(5R)\) 后没有约去一个 \(M\)。
\item
  同密度时误认为 \(R\propto M\)，从而得到错误的质量标度。
\end{itemize}

\subsection{八、课堂讲解建议}\label{ux516bux8bfeux5802ux8bb2ux89e3ux5efaux8bae-5}

先用``逐层搬运球壳''的图像推导
\(3/5\)，再进行数值代入。第二问算完后应与太阳年龄直接比较，让学生看到引力供能不足的数量级证据。第三问先做符号比例，再代数值，避免大数运算掩盖
\(M^{2/3}\) 的物理结构。

数值估算题的课堂重点应放在结构而不是按计算器。先在黑板左侧保留符号式
\(E_g\sim GM^2/R\)、\(P\sim R^2\sigma T^4\)、\(t=E/P\)，右侧再逐项代入指数。让学生先预估数量级，再计算
\(3/5\)、\(4\pi\) 等系数，能够及时发现指数错位。

推导 \(3/5\)
时用同心薄壳装配图说明每一薄壳只与内部已有质量相互作用。第二问算出开尔文---亥姆霍兹时间后，直接与太阳约
\(10^{10}\)
年的年龄比较；这个比较比单独报告秒数更有物理意义。第三问强调题设采用``全部转化为内能''的简化，不能擅自加入课外模型系数。

质量标度建议分两步讲：同密度给出 \(R\propto M^{1/3}\)，再由
\(T\propto GM/R\) 得到
\(T\propto M^{2/3}\)。先完成幂律再代点火温度，可使学生看到估算问题中比例关系往往比绝对数值更可迁移。

\subsection{九、变式训练}\label{ux4e5dux53d8ux5f0fux8badux7ec3-5}

\begin{enumerate}
\def\labelenumi{\arabic{enumi}.}
\tightlist
\item
  若恒星密度不变但比热容为太阳的两倍，求点火质量比例如何改变。
\item
  对密度分布不均匀、引力自能写成 \(-aGM^2/R\) 的恒星，讨论系数 \(a\)
  对结果的影响。
\end{enumerate}

第一题中温升变为原来的一半，要达到同一点火温度需提高 \(M/R\)；结合
\(R\propto M^{1/3}\)
可求新的质量幂律。第二题训练学生保留结构系数，最终应得到临界质量随
\((c/a)^{3/2}\) 缩放，并说明质量更向中心集中的恒星通常有更大的 \(a\)。

还可把``当前辐射功率保持不变''改为主序星经验关系
\(L\propto M^{3.5}\)，讨论不同质量恒星的引力收缩时间如何标度。该题训练把能量标度和功率标度组合，而不要求追求精确天体模型。

\textbf{变式解答。} 同密度条件下
\(T\propto M^{2/3}/c\)。若比热加倍，达到同一点火温度的质量下限变为原值的
\(2^{3/2}\) 倍，即约 \(0.207M\)。若只保留收缩能的比例 \(f\)，则
\(T=fT_{\rm full}\)，相应临界质量为 \(f^{-3/2}\)
倍。以上是题设均匀密度和常比热模型的结果，不把这个数值直接视为真实恒星的精确质量界限。

\newpage

\section{第 7 题}\label{ux7b2c-7-ux9898}

\(1\,\mathrm{mol}\) 单原子分子理想气体在温度---熵图上经历圆形循环。令

\[x=\frac{T}{T_0},\qquad y=\frac{S}{S_0},\qquad S_0=2R,\]

循环满足

\[(x-2)^2+(y-2)^2=1.\]

规定气体处于 \((T_0,V_0,p_0)\) 状态时熵为 \(S_0\)。

\begin{enumerate}
\def\labelenumi{\arabic{enumi}.}
\tightlist
\item
  若循环为正循环，即气体在一个循环中对外做正功，判断图中的循环方向。
\item
  求循环效率。
\item
  求循环中的最大体积比 \(V_m/V_0\)。
\end{enumerate}

\begin{figure}[H]
\centering
\includegraphics[width=0.569\linewidth,height=\textheight,keepaspectratio,alt={图 7-1　温度---熵平面上的圆形循环。}]{../figures/question-07.pdf}
\caption*{图 7-1　温度---熵平面上的圆形循环。}
\end{figure}

\subsection{一、题型定位}\label{ux4e00ux9898ux578bux5b9aux4f4d-6}

在温度---熵图上沿路径积分
\(T\,dS\)，求吸热、放热与循环效率。最后将理想气体熵公式代入圆周约束，求体积极值。

\subsection{二、知识点}\label{ux4e8cux77e5ux8bc6ux70b9-6}

\begin{itemize}
\tightlist
\item
  可逆过程 \(\delta Q=T\,\mathrm dS\)。
\item
  循环过程 \(\Delta U=0\)，故 \(W=Q_{\mathrm{net}}\)。
\item
  \(T\!-!S\) 图的有向面积和吸热区间。
\item
  单原子理想气体熵变 \[S-S_0=R\ln(V/V_0)+\frac32R\ln(T/T_0).\]
\item
  拉格朗日乘数法或圆参数法求最大体积。
\end{itemize}

\subsection{三、物理图像}\label{ux4e09ux7269ux7406ux56feux50cf-6}

在 \(T\) 为横轴、\(S\)
为纵轴的图上，沿逆时针方向绕行时，右半圆熵增加并吸热，左半圆熵减少并放热。净吸热量等于圆的有向面积，因此逆时针循环对外做正功。

体积并不只随熵增大而增大：在熵相同的情况下，温度越高，体积越小。因此最大体积点位于圆的左上象限，而不是最高点或最左点。

\subsection{四、解题入口}\label{ux56dbux89e3ux9898ux5165ux53e3-6}

引入无量纲变量

\[x=\frac{T}{T_0},\qquad y=\frac{S}{S_0},\]

圆方程为

\[(x-2)^2+(y-2)^2=1. \tag{1}\]

计算效率时，可取逆时针参数

\[x=2+\cos\phi,\qquad y=2+\sin\phi.\]

计算最大体积时，先把 \(\ln(V/V_0)\) 写成 \(x,y\)
的函数，再在圆上求极值。

\subsection{五、解析}\label{ux4e94ux89e3ux6790-6}

图 7-2 在无量纲 \(x=T/T_0\)、\(y=S/S_0\)
平面中画出圆形循环。逆时针方向使 \(\oint T\,dS\) 为正；右半圆满足
\(dS>0\)，对应吸热；黑点标出由熵状态式求得的最大体积位置。第（1）、（2）、（3）步分别使用这三个几何信息。

\begin{figure}[H]
\centering
\includegraphics[width=0.667\linewidth,height=\textheight,keepaspectratio,alt={图 7-2　逆时针正循环、吸热半圆与最大体积点。}]{../figures/solution-07.pdf}
\caption*{图 7-2　逆时针正循环、吸热半圆与最大体积点。}
\end{figure}

\textbf{（1）循环方向}

在图 7-2 中，循环功 \(W=\oint T\,dS\) 等于有向面积乘尺度
\(T_0S_0\)。逆时针沿右侧向上、左侧向下，积分为正，所以这是正循环方向。

可逆过程中

\[\delta Q=T\,\mathrm dS. \tag{2}\]

一个循环后内能回到原值，\(\Delta U=0\)，所以

\[W=Q_{\mathrm{net}}=\oint T\,\mathrm dS. \tag{3}\]

取逆时针参数

\[T=T_0(2+\cos\phi),\qquad
S=S_0(2+\sin\phi),\]

其中 \(0\leq\phi\leq2\pi\)。于是

\[\mathrm dS=S_0\cos\phi\,\mathrm d\phi.\]

参数点的运动方向也应核对：\(\phi=0\) 时位于圆的最右端，且

\[\frac{\mathrm dT}{\mathrm d\phi}=0,\qquad
\frac{\mathrm dS}{\mathrm d\phi}=S_0>0,\]

所以随 \(\phi\)
增大，状态点从最右端向上运动，确为逆时针。这样可在积分之前固定正负号，而不是由最后算出的功为正再倒推方向。

净功为

\[\begin{aligned}
W
&=S_0T_0\int_0^{2\pi}(2+\cos\phi)\cos\phi\,\mathrm d\phi\\
&=S_0T_0\left[2\int_0^{2\pi}\cos\phi\,\mathrm d\phi
+\int_0^{2\pi}\cos^2\phi\,\mathrm d\phi\right]\\
&=\pi S_0T_0>0.
\end{aligned} \tag{4}\]

故正循环方向为

\[\boxed{\text{逆时针}}. \tag{5}\]

\textbf{（2）循环效率}

图 7-2 的右半圆沿逆时针方向有 \(dS>0\)，故
\(\delta Q=T\,dS>0\)，是吸热段；左半圆是放热段。效率分子由整圆有向面积给出，分母则必须单独积分右半圆，不能用半个圆面积代替热量。

逆时针循环中，熵增加的右半圆为吸热段，对应

\[-\frac\pi2\leq\phi\leq\frac\pi2.\]

吸热量

\[\begin{aligned}
Q_{\mathrm{in}}
&=S_0T_0\int_{-\pi/2}^{\pi/2}
(2+\cos\phi)\cos\phi\,\mathrm d\phi\\
&=S_0T_0\left[
2\int_{-\pi/2}^{\pi/2}\cos\phi\,\mathrm d\phi
+\int_{-\pi/2}^{\pi/2}\cos^2\phi\,\mathrm d\phi
\right]\\
&=S_0T_0\left(4+\frac\pi2\right). 
\end{aligned} \tag{6}\]

其中两个定积分分别为

\[\int_{-\pi/2}^{\pi/2}\cos\phi\,\mathrm d\phi
=\left.\sin\phi\right|_{-\pi/2}^{\pi/2}=2,\]

以及

\[\begin{aligned}
\int_{-\pi/2}^{\pi/2}\cos^2\phi\,\mathrm d\phi
&=\frac12\int_{-\pi/2}^{\pi/2}(1+\cos2\phi)\,\mathrm d\phi\\
&=\frac\pi2.
\end{aligned}\]

这里使用了降幂公式
\(\cos^2\phi=(1+\cos2\phi)/2\)。把两项分别算清楚，可以直接看出 \(4\)
来自平均温度部分，\(\pi/2\) 来自圆周温度起伏部分。

由式（4），循环做功为

\[W=\pi S_0T_0. \tag{7}\]

所以效率

\[\boxed{
\eta=\frac{W}{Q_{\mathrm{in}}}
=\frac{\pi}{4+\pi/2}
=\frac{2\pi}{8+\pi}
}. \tag{8}\]

\textbf{（3）最大体积}

黑点不是圆的最右点或最高点，因为体积同时依赖 \(T\) 与 \(S\)。先把
\(\ln(V/V_0)\) 写成图中坐标 \(x,y\)
的函数，再在圆周约束下沿切向令方向导数为零，才能找到最大体积点。

对 \(1\,\mathrm{mol}\) 单原子理想气体，取 \((T_0,V_0,S_0)\) 为参考状态：

\[S-S_0=R\ln\frac{V}{V_0}+\frac32R\ln\frac{T}{T_0}. \tag{9}\]

因为 \(S_0=2R\)、\(S=yS_0\)、\(T=xT_0\)，式（9）变为

\[2R(y-1)=R\ln\frac{V}{V_0}+\frac32R\ln x.\]

约去 \(R\)：

\[\ln\frac{V}{V_0}=2(y-1)-\frac32\ln x. \tag{10}\]

令

\[F(x,y)=2(y-1)-\frac32\ln x.\]

在约束式（1）上求 \(F\) 的最大值。引入拉格朗日乘子 \(\Lambda\)：

\[\nabla F=\Lambda\nabla[(x-2)^2+(y-2)^2].\]

分量方程为

\[-\frac{3}{2x}=2\Lambda(x-2),\qquad
2=2\Lambda(y-2). \tag{11}\]

由第二式

\[\Lambda=\frac1{y-2}.\]

代回第一式，得到

\[4x(x-2)+3(y-2)=0. \tag{12}\]

将式（12）与圆方程联立，选择左上象限的极大值点，数值解为

\[x_m\approx1.56865,\qquad y_m\approx2.90218. \tag{13}\]

若希望把数值求解的对象写成单变量方程，可先由式（12）表示

\[y-2=-\frac{4x(x-2)}{3}.\]

代入圆约束：

\[ (x-2)^2+\frac{16x^2(x-2)^2}{9}=1.\]

提取公因子后得到

\[ (x-2)^2(9+16x^2)=9. \tag{13a}\]

在圆上 \(1\leq x\leq3\)。逐一求式（13a）在该区间内的实根，再由

\[y=2-\frac{4x(x-2)}{3}\]

恢复
\(y\)。位于左上象限的根给出式（13）；其余驻点要么位于右下象限，要么对应体积极小值。判别极大值最稳妥的方法，是把每个候选点代回

\[F(x,y)=2(y-1)-\frac32\ln x\]

直接比较，而不是只凭象限选择。

把它代入式（10）：

\[\begin{aligned}
\ln\frac{V_m}{V_0}
&=2(2.902-1)-\frac32\ln(1.568)\\
&\approx3.129.
\end{aligned}\]

因此

\[\boxed{\frac{V_m}{V_0}\approx e^{3.129}\approx22.85}. \tag{14}\]

在 \(T\)---\(S\)
图上必须从热力学基本关系开始确定方向和面积含义。可逆过程中
\(\delta Q=T\,\mathrm dS\)，所以一周吸收的净热量为
\(\oint T\,\mathrm dS\)；完整循环有
\(\Delta U=0\)，故净功也等于该积分。由于本题把 \(T\) 画在横轴、\(S\)
画在纵轴，\(\oint T\,\mathrm dS\) 对逆时针圆周为正，这与常见
\(p\)---\(V\) 图``顺时针做正功''的视觉习惯相反。

计算效率时，分母不是整个圆周上 \(T\,\mathrm dS\) 的绝对值之和，而是
\(\mathrm dS>0\)
的吸热半圆上的积分。参数化时应明确角参数的起点和增加方向，把
\(T=T_0(2+\cos\phi)\)、\(S=S_0(2+\sin\phi)\) 代入，再只在 \(\cos\phi>0\)
的弧段积分。净功等于圆在物理坐标中的面积
\(\pi T_0S_0\)，可作为积分结果的独立检查。

最大体积不能从图上直接取最大熵点。单原子理想气体满足

\[S-S_0=R\ln\frac{V}{V_0}+\frac32R\ln\frac{T}{T_0},\]

因此应最大化

\[2(y-1)-\frac32\ln x.\]

拉格朗日乘子方程给出的驻点可能有多个，必须代回圆约束并比较候选点；最大值位于左上象限，体现``较大熵''和``较低温度''共同增大体积。最后代回原状态方程而不是只报告数值坐标。

效率表达式也可做上下界检查。由于循环中既有吸热又有放热，必有
\(0<\eta<1\)。本题结果

\[\eta=\frac{2\pi}{8+\pi}\approx0.56\]

满足这一要求。若误把整个圆周上的 \(|T\,\mathrm dS|\)
当作吸热量，分母会过大；若把净功圆面积直接当作吸热量，则会错误得到效率1。最大体积点求出后还应计算圆约束残差，确认数值解没有因舍入离开轨迹。

温度---熵图上的方向和面积都来自 \(\delta Q=T\,\mathrm dS\)：

\begin{enumerate}
\def\labelenumi{\arabic{enumi}.}
\tightlist
\item
  先按参数角写 \(T(\phi),S(\phi)\) 与
  \(\mathrm dS\)，再积分一周；有向面积的符号决定正循环方向，不能凭普通
  \(p\)---\(V\) 图经验判断。
\item
  吸热段由 \(\mathrm dS>0\) 且 \(T>0\)
  确定，积分上下限应覆盖右半圆。净功是闭合积分，效率分母则只取正热量，两者不能混用。
\item
  最大体积从熵公式解出
  \(\ln(V/V_0)\)，代入圆周参数后变成单变量极值；对该函数求导并检查两个驻点及端点。
\item
  最后把最优参数代回圆方程和熵公式，检查 \(V_m/V_0>1\) 且效率位于 \(0\)
  与 \(1\) 之间。
\end{enumerate}

\textbf{热量的路径积分：} 在 \(T\)---\(S\)
图上，循环功和吸热量的积分域不同：

\[\begin{aligned}
W_{\rm net}&=\oint T\,\mathrm dS,\\
Q_{\rm in}&=\int_{\mathrm dS>0}T\,\mathrm dS,\\
\eta&=\frac{\oint T\,\mathrm dS}
{\int_{\mathrm dS>0}T\,\mathrm dS}.
\end{aligned}\]

若用参数 \(\phi\) 描述圆周，必须把
\(\mathrm dS=(\mathrm dS/\mathrm d\phi)\mathrm d\phi\) 一并代入；只积分
\(T(\phi)\) 会漏掉雅可比因子。

\subsection{六、关键洞察}\label{ux516dux5173ux952eux6d1eux5bdf-6}

\(T\!-!S\) 图与通常的 \(p\!-!V\) 图不同：这里 \(T\) 是横轴、\(S\)
是纵轴，净功对应
\(\oint T\,\mathrm dS\)，逆时针为正。最大体积是熵增大和温度降低共同作用的结果，必须优化组合量
\(2(y-1)-\frac32\ln x\)。

\subsection{七、易错点}\label{ux4e03ux6613ux9519ux70b9-6}

\begin{itemize}
\tightlist
\item
  把 \(T\!-!S\) 图的正循环方向照搬成 \(p\!-!V\) 图的顺时针。
\item
  把吸热量误写成圆的上半部分面积；吸热条件是
  \(\mathrm dS>0\)，对应右半圆。
\item
  效率分母用净吸热而不是总吸热 \(Q_{\mathrm{in}}\)。
\item
  熵公式中遗漏单原子气体的 \(3R/2\)。
\item
  认为最大熵点就是最大体积点。
\item
  数值极值求出后没有检查它位于圆上且确为最大值。
\end{itemize}

\subsection{八、课堂讲解建议}\label{ux516bux8bfeux5802ux8bb2ux89e3ux5efaux8bae-6}

先让学生沿圆周用箭头判断哪些弧段
\(\mathrm dS>0\)，再积分。效率计算可把净功解释为圆面积，把吸热量解释为右半圆的曲线积分。第三问先定性猜测极值位于左上象限，再用拉格朗日方程验证。

开场先对比 \(p\)---\(V\) 图和 \(T\)---\(S\)
图的面积含义。让学生亲自沿逆时针和顺时针各走一圈，用
\(\delta Q=T\,\mathrm dS\)
判断净热量符号，避免套用``顺时针做正功''的旧记忆。随后在圆上标出
\(\mathrm dS>0\) 的右半圆，明确只有这一段进入吸热量分母。

积分前先约定参数角的零点和增加方向，并把 \(\mathrm dS\)
写出来。净功既可直接积分，也可用圆面积求得；两种方法相互核验。第三问不宜直接进入拉格朗日乘子，应先由熵公式说明增大
\(S\)、减小 \(T\) 都可能增大 \(V\)，因此极值应在左上象限附近。

求出驻点后让学生把坐标代回圆方程，并比较所有候选点。最后再由
\(\ln(V/V_0)\)
还原体积比。整个过程可以训练``先定性定位---再建立目标函数---最后筛选驻点''的约束极值流程。

\subsection{九、变式训练}\label{ux4e5dux53d8ux5f0fux8badux7ec3-6}

\begin{enumerate}
\def\labelenumi{\arabic{enumi}.}
\tightlist
\item
  把圆的半径改为无量纲参数 \(a<2\)，求净功和吸热量的一般表达式。
\item
  把工质改为 \(1\,\mathrm{mol}\) 双原子理想气体，重新计算最大体积点。
\end{enumerate}

第一题训练参数化积分，净功应随圆面积按 \(a^2\)
缩放，而吸热量还含圆心温度与半径共同产生的项。第二题只改变熵状态方程中
\(\ln T\)
的系数，因此循环效率不变，但最大体积点会移动；这能区分``路径吸放热''与``工质状态方程''分别影响哪些结果。

还可把圆改为椭圆
\((x-2)^2/a^2+(y-2)^2/b^2=1\)，要求只求净功和循环方向。学生应直接利用椭圆面积而不做冗长积分，并说明只要曲线不跨入
\(T<0\) 区域，方向判据仍由 \(\oint T\,\mathrm dS\) 决定。

\textbf{变式解答：圆心温度改变。} 若改成
\(x=a+\cos\phi\)、\(y=b+\sin\phi\)，\(a>1\)
保证正温度，正循环仍为逆时针。净功 \(W=\pi T_0S_0\)，吸热量
\(Q_{\rm in}=T_0S_0(2a+\pi/2)\)，故

\[\eta=\frac{\pi}{2a+\pi/2}.\]

圆心熵 \(b\) 不进入效率；它通过状态方程改变绝对体积范围。取 \(a=2\)
恢复正文结果。

\newpage

\section{第 8 题}\label{ux7b2c-8-ux9898}

假设存在磁荷为 \(g_m\) 的磁单极子，它在距自身 \(r\) 处产生的磁场为

\[\boldsymbol B=\frac{g_m\boldsymbol r}{4\pi r^3}.\]

磁单极子固定在坐标原点。电量为 \(q\) 的带电粒子从 \(z=-\infty\) 处沿直线
\(x=b\) 飞来，远处速度为
\(\boldsymbol v_0=v\hat{\boldsymbol z}\)。认为磁场造成的偏转为小量，可用未偏转的直线轨迹计算最低阶效应。

\begin{enumerate}
\def\labelenumi{\arabic{enumi}.}
\tightlist
\item
  求粒子通过磁场后的动量变化 \(\Delta\boldsymbol p\)，保留最低阶非零项。
\item
  假设偏转过程中粒子机械角动量的变化量满足玻尔式量子化条件。若由此推出的电荷最小单位应等于基本电荷
  \(e\)，求磁荷 \(g_m\)，用 \(e\) 和普朗克常量 \(h\) 表示。
\end{enumerate}

\begin{figure}[H]
\centering
\includegraphics[width=0.756\linewidth,height=\textheight,keepaspectratio,alt={图 8-1　带电粒子以碰撞参数 b 掠过磁单极子。}]{../figures/question-08.pdf}
\caption*{图 8-1　带电粒子以碰撞参数 \(b\) 掠过磁单极子。}
\end{figure}

\subsection{一、题型定位}\label{ux4e00ux9898ux578bux5b9aux4f4d-7}

小偏转时沿未扰动直线积分洛伦兹力，先求横向冲量，再算机械角动量变化。量子化关系按题设使用，且结果须检查小偏转条件。

\subsection{二、知识点}\label{ux4e8cux77e5ux8bc6ux70b9-7}

\begin{itemize}
\tightlist
\item
  磁单极子径向磁场 \(\boldsymbol B=g_m\boldsymbol r/(4\pi r^3)\)。
\item
  洛伦兹力 \(\boldsymbol F=q\boldsymbol v\times\boldsymbol B\)。
\item
  小偏转近似中的直线轨迹积分。
\item
  标准积分
  \[\int_{-\infty}^{\infty}\frac{\mathrm dz}{(b^2+z^2)^{3/2}}=\frac{2}{b^2}.\]
\item
  冲量与角动量变化、\(\hbar=h/(2\pi)\)。
\end{itemize}

\subsection{三、物理图像}\label{ux4e09ux7269ux7406ux56feux50cf-7}

未偏转轨迹位于 \(x=b\)、\(y=0\) 的直线上。磁场既有 \(x\) 分量也有 \(z\)
分量，但 \(z\) 分量与粒子速度平行，不产生磁力；只有 \(B_x\) 产生沿 \(y\)
方向的力。粒子从 \(z=-\infty\) 到 \(+\infty\) 的整个过程中，\(B_x\)
同号，因此横向冲量不会相互抵消。

\subsection{四、解题入口}\label{ux56dbux89e3ux9898ux5165ux53e3-7}

在最低阶近似中取

\[\boldsymbol r=b\hat{\boldsymbol x}+z\hat{\boldsymbol z},\qquad
\boldsymbol v=v\hat{\boldsymbol z},\qquad
r=\sqrt{b^2+z^2}.\]

先直接计算叉积的方向，再用 \(\mathrm dt=\mathrm dz/v\)
消去速度。角动量变化取关于原点的机械角动量，最低阶为 \(b\Delta p_y\)。

\subsection{五、解析}\label{ux4e94ux89e3ux6790-7}

图 8-2 取电荷沿 \(+z\) 方向以速度 \(v\)
经过磁单极子，冲量近似中使用未偏转直线 \(x=b\)。径向磁场在轨迹点分解为
\(x,z\) 分量；其中 \(B_x\) 与 \(v\hat{\bm z}\)
叉乘产生垂直纸面的横向力，积分得到 \(\Delta p_y\)，再以力臂 \(b\)
求角动量变化。

\begin{figure}[H]
\centering
\includegraphics[width=0.875\linewidth,height=\textheight,keepaspectratio,alt={图 8-2　未偏转直线轨迹上的磁场分量、洛伦兹力与横向冲量。}]{../figures/solution-08.pdf}
\caption*{图 8-2　未偏转直线轨迹上的磁场分量、洛伦兹力与横向冲量。}
\end{figure}

\textbf{（1）动量变化}

按图 8-2，\(\bm v\parallel\hat{\bm z}\)，磁场的 \(z\)
分量与速度平行而不产生洛伦兹力；只有指向横向的 \(B_x\) 产生 \(y\)
向力。圆点表示 \(+y\) 方向出纸面，若 \(qg_m\) 变号则冲量方向反转。

沿未偏转轨迹，磁场为

\[\boldsymbol B(z)
=\frac{g_m}{4\pi(b^2+z^2)^{3/2}}
\left(b\hat{\boldsymbol x}+z\hat{\boldsymbol z}\right). \tag{1}\]

洛伦兹力

\[\boldsymbol F=q\boldsymbol v\times\boldsymbol B.\]

代入 \(\boldsymbol v=v\hat{\boldsymbol z}\)：

\[\begin{aligned}
\boldsymbol F
&=qv\hat{\boldsymbol z}\times
\frac{g_m(b\hat{\boldsymbol x}+z\hat{\boldsymbol z})}
{4\pi(b^2+z^2)^{3/2}}\\
&=\frac{qg_mvb}{4\pi(b^2+z^2)^{3/2}}
(\hat{\boldsymbol z}\times\hat{\boldsymbol x})\\
&=\frac{qg_mvb}{4\pi(b^2+z^2)^{3/2}}
\hat{\boldsymbol y}.
\end{aligned} \tag{2}\]

这里使用了
\(\hat{\boldsymbol z}\times\hat{\boldsymbol x}=\hat{\boldsymbol y}\)，而
\(\hat{\boldsymbol z}\times\hat{\boldsymbol z}=0\)。

动量变化为冲量：

\[\Delta\boldsymbol p=\int_{-\infty}^{\infty}\boldsymbol F\,\mathrm dt.\]

由于 \(z=vt\)，有 \(\mathrm dt=\mathrm dz/v\)，所以

\[\begin{aligned}
\Delta\boldsymbol p
&=\frac{qg_mvb}{4\pi}\hat{\boldsymbol y}
\int_{-\infty}^{\infty}
\frac{\mathrm dz/v}{(b^2+z^2)^{3/2}}\\
&=\frac{qg_mb}{4\pi}\hat{\boldsymbol y}
\int_{-\infty}^{\infty}
\frac{\mathrm dz}{(b^2+z^2)^{3/2}}.
\end{aligned} \tag{3}\]

令 \(z=b\tan u\)，则

\[\mathrm dz=b\sec^2u\,\mathrm du,\qquad
(b^2+z^2)^{3/2}=b^3\sec^3u.\]

因此

\[\begin{aligned}
\int_{-\infty}^{\infty}
\frac{\mathrm dz}{(b^2+z^2)^{3/2}}
&=\frac1{b^2}\int_{-\pi/2}^{\pi/2}\cos u\,\mathrm du\\
&=\frac2{b^2}.
\end{aligned} \tag{4}\]

也可以用一个便于记忆的原函数直接检查。因为

\[\frac{\mathrm d}{\mathrm dz}
\left(\frac{z}{\sqrt{b^2+z^2}}\right)
=\frac{b^2}{(b^2+z^2)^{3/2}},\]

所以

\[\int\frac{\mathrm dz}{(b^2+z^2)^{3/2}}
=\frac{z}{b^2\sqrt{b^2+z^2}}+C.\]

在上下限代入时，

\[\lim_{z\to+\infty}\frac{z}{\sqrt{b^2+z^2}}=1,\]

\[\lim_{z\to-\infty}\frac{z}{\sqrt{b^2+z^2}}=-1,\]

故定积分仍为
\(2/b^2\)。三角代换和原函数两种算法互相核验，可排除把答案误写成
\(\pi/b^2\) 的常见错误。

代回式（3）：

\[\boxed{
\Delta\boldsymbol p
=\frac{qg_m}{2\pi b}\hat{\boldsymbol y}
}. \tag{5}\]

若
\(qg_m<0\)，方向反向；式（5）的符号已经自动包含这一点。纵向动量变化要到更高阶才出现。

\textbf{（2）磁荷量子化条件}

图 8-2 中横向冲量的作用线距原点约为 \(b\)，因此机械角动量改变量是
\(b\Delta p_y\)。代入上一问后冲量中的 \(1/b\)
恰好抵消，量子化条件由此与具体掠射距离无关。

最低阶横向冲量作用在距原点约为 \(b\)
的轨迹上，因此机械角动量变化量的大小为

\[\Delta L=b|\Delta p_y|
=\frac{|qg_m|}{2\pi}. \tag{6}\]

若要同时核对方向，可直接积分力矩。由

\[\boldsymbol r=b\hat{\boldsymbol x}+z\hat{\boldsymbol z},\]

\[\boldsymbol F=F_y\hat{\boldsymbol y},\]

得到

\[\begin{aligned}
\boldsymbol\tau
&=\boldsymbol r\times\boldsymbol F\\
&=bF_y\hat{\boldsymbol z}-zF_y\hat{\boldsymbol x}.
\end{aligned}\]

\(F_y(z)\) 是 \(z\) 的偶函数，因此 \(-zF_y\) 是奇函数，从 \(-\infty\) 到
\(+\infty\) 的积分为零：

\[\Delta L_x=-\int_{-\infty}^{\infty}zF_y\,\mathrm dt=0.\]

而

\[\Delta L_z=b\int_{-\infty}^{\infty}F_y\,\mathrm dt
=b\Delta p_y
=\frac{qg_m}{2\pi}.\]

所以式（6）不仅给出大小，也说明在 \(qg_m>0\)、\(b>0\)
的约定下角动量增量沿 \(+z\) 方向。

按题目给出的玻尔式量子化假设，令

\[\Delta L=n\hbar=n\frac{h}{2\pi},\qquad n\in\mathbb Z. \tag{7}\]

比较式（6）、（7），得到

\[|qg_m|=nh. \tag{8}\]

若固定磁荷 \(g_m\) 后允许的最小正电荷应为实验基本电荷
\(e\)，取最小非零量子数 \(n=1\)，则

\[eg_m=h.\]

因此

\[\boxed{g_m=\frac{h}{e}}. \tag{9}\]

这里的结果遵从题目给定的磁场定义和简化量子化假设；不同单位制中的磁荷定义可能把常数因子吸收到
\(g_m\) 中。

冲量方向最好先用坐标分量判断。未偏转轨迹上

\[\boldsymbol r=b\hat{\boldsymbol x}+vt\hat{\boldsymbol z},\qquad
\boldsymbol v=v\hat{\boldsymbol z},\]

所以 \(\boldsymbol v\times\boldsymbol B\) 中只有磁场的 \(x\)
分量贡献净横向力，方向为 \(+\hat{\boldsymbol y}\)；含 \(z\)
分量的叉积为零。于是

\[F_y=\frac{qg_mvb}{4\pi(b^2+v^2t^2)^{3/2}}.\]

令 \(u=vt/b\) 后，\(\mathrm dt=(b/v)\mathrm du\)，所有速度因子和一个
\(b\) 抵消，剩下的标准积分为

\[\int_{-\infty}^{\infty}\frac{\mathrm du}{(1+u^2)^{3/2}}=2.\]

这一步同时说明冲量与 \(v\) 无关、与 \(b\) 成反比。

角动量变化取 \(z\) 轴附近的入射几何最方便。最低阶中横向冲量作用的力臂为
\(b\)，因此

\[|\Delta L|=b|\Delta p_y|.\]

其方向可由 \(\boldsymbol r\times\Delta\boldsymbol p\) 再检查一次。\(b\)
在乘积中消失后，量子化条件不再依赖散射轨迹的具体尺度。比较

\[\frac{qg_m}{2\pi}=n\hbar,\qquad \hbar=\frac{h}{2\pi}\]

可见两个 \(2\pi\) 正好消去，得到题目所用约定下的
\(qg_m=nh\)。若换用其他磁荷单位制，必须先换定义，不能直接搬用本题结论。

磁力始终与瞬时速度垂直，因此不改变粒子动能。最低阶近似中只产生横向动量，一阶纵向动量变化为零；为了保持总动量大小不变，纵向修正从偏转角的二阶才出现。冲量方向随
\(qg_m\) 的符号反转，而量子化条件通常写绝对值。若
\(b<0\)，几何力臂和横向冲量方向同时改变，\(\Delta L\)
的物理结论仍应按统一的有向坐标重新判断。

磁单极模型中的冲量应从未偏转轨迹逐分量积分：

\begin{enumerate}
\def\labelenumi{\arabic{enumi}.}
\tightlist
\item
  写 \(\boldsymbol r=b\hat x+vt\hat z\) 和磁场
  \(\boldsymbol B\propto\boldsymbol r/r^3\)，再计算
  \(q\boldsymbol v\times\boldsymbol B\)；纵向分量在最低阶为零，横向方向由叉乘确定。
\item
  将 \(u=vt/b\) 代入时间积分，完整写出 \(\mathrm dt=(b/v)\mathrm du\)
  与分母 \((1+u^2)^{3/2}\)，积分上下限由 \(t=\pm\infty\) 变为
  \(u=\pm\infty\)。
\item
  角动量变化用冲量的力臂 \(b\)
  相乘，随后按题设的量子化条件解磁荷；不能把 \(h\) 与 \(\hbar\) 混用。
\item
  检查结果与粒子速率无关，这是力随作用时间的补偿结果；若积分后仍含
  \(v\)，应回查变量替换。
\end{enumerate}

\textbf{冲量积分：} 未偏转近似下最常用的积分可以完整写成

\[\begin{aligned}
I&=\int_{-\infty}^{+\infty}
\frac{v\,\mathrm dt}{(b^2+v^2t^2)^{3/2}},\\
u&=\frac{vt}{b},\qquad \mathrm dt=\frac b v\,\mathrm du,\\
I&=\frac1{b^2}\int_{-\infty}^{+\infty}
\frac{\mathrm du}{(1+u^2)^{3/2}}
=\frac{2}{b^2}.
\end{aligned}\]

横向冲量再乘相应的场强系数，角动量改变量为
\(\Delta L=b\Delta p_\perp\)；因此最终的 \(v\) 必须完全消去。

\subsection{六、关键洞察}\label{ux516dux5173ux952eux6d1eux5bdf-7}

动量冲量虽与碰撞参数成 \(1/b\)，角动量还要乘力臂 \(b\)，两者正好抵消，使
\(\Delta L=qg_m/(2\pi)\)
与具体轨迹尺度无关。这种拓扑式的参数独立性正是从磁单极子推出电荷量子化的核心。

\subsection{七、易错点}\label{ux4e03ux6613ux9519ux70b9-7}

\begin{itemize}
\tightlist
\item
  把 \(\boldsymbol v\times\boldsymbol B\) 的方向写成
  \(-\hat{\boldsymbol y}\)。
\item
  把磁场的 \(B_z\) 也计入磁力；它与速度平行，不产生力。
\item
  积分时同时保留 \(v\) 和 \(\mathrm dt\)，忘记用
  \(\mathrm dz=v\,\mathrm dt\) 消去速度。
\item
  标准积分误算成 \(\pi/b^2\)。
\item
  角动量变化误写成 \(\Delta p_y/b\) 而不是 \(b\Delta p_y\)。
\item
  把 \(h\) 与 \(\hbar\) 混淆，遗漏 \(2\pi\) 的抵消。
\end{itemize}

\subsection{八、课堂讲解建议}\label{ux516bux8bfeux5802ux8bb2ux89e3ux5efaux8bae-7}

先让学生只做向量叉积，明确为什么力始终沿同一个横向方向。积分完成后追问结果为什么与粒子速度无关。第二问应把
\(1/b\) 与力臂 \(b\) 写在同一行，让学生直观看到角动量变化的普适性。

建议把轨迹写成
\(\boldsymbol r(t)=b\hat{\boldsymbol x}+vt\hat{\boldsymbol z}\)，然后分解磁场的
\(x\)、\(z\) 分量。让学生逐项做
\(\hat{\boldsymbol z}\times\hat{\boldsymbol x}\) 和
\(\hat{\boldsymbol z}\times\hat{\boldsymbol z}\)，先确定只有 \(y\)
向冲量，再进入积分。这样能够避免只算大小、最后猜方向。

无量纲代换 \(u=vt/b\)
应完整写出微分变换。积分后用三项检查巩固结果：冲量单位正确；\(b\)
越大偏转越小；磁力不做功，因此最低阶速度大小不变。第二问把
\(b\Delta p_y\) 写成一行，突出 \(b\) 消去后角动量变化与轨道尺度无关。

最后说明量子化步骤依赖题目给定的磁荷定义。让学生先写
\(\Delta L=n\hbar\)，再显式代入
\(\hbar=h/(2\pi)\)；如果跳过这一行，最容易在答案中多出或少掉 \(2\pi\)。

\subsection{九、变式训练}\label{ux4e5dux53d8ux5f0fux8badux7ec3-7}

\begin{enumerate}
\def\labelenumi{\arabic{enumi}.}
\tightlist
\item
  若入射直线改为 \(y=b\)、速度仍沿 \(+z\)，求冲量方向。
\item
  保留小偏转的下一阶效应，判断纵向动量变化的数量级，并说明磁力为何不做功。
\end{enumerate}

第一题训练右手定则和坐标轮换：先写新的位置矢量，再做叉积，冲量大小不变而方向旋转。第二题训练微扰阶数；横向速度是一阶小量，保持速率不变要求纵向速度修正从二阶开始，因此纵向动量变化不是本题最低阶非零结果。

可进一步让磁单极子缓慢移动，要求判断实验室系中是否仍可只用静磁场洛伦兹力。该讨论用于提醒学生：本题的简洁结论依赖单极子固定以及未偏转直线近似，改变源的运动状态后还需考虑相应的电场或相对论变换。

\textbf{变式解答：最低阶的适用范围。}
偏转角与直线冲量近似的适用条件分别为

\[\chi\approx\frac{|qg_m|}{2\pi bmv},\qquad |qg_m|\ll2\pi bmv.\]

磁力不做功，动量大小保持 \(mv\)；若已算出一阶
\(\Delta p_y\)，则纵向最低非零修正为

\[\Delta p_z=-\frac{(\Delta p_y)^2}{2mv}+O(\chi^4mv).\]

它是二阶项，说明``纵向无变化''只在一阶成立。题目给的机械角动量变化量量子化是假设，本解据此得到
\(g_m=h/e\)，不将它等同于其他磁荷定义或完整量子理论的量子化推导。

\end{document}
